\documentclass[11pt]{article}
\usepackage[a4paper,margin=27mm]{geometry}
\usepackage{amsmath,amssymb,amsthm,mathtools}
\usepackage[T1]{fontenc}
\usepackage{lmodern}
\usepackage{xcolor}
\newtheorem{theorem}{Theorem}
\newtheorem{lemma}[theorem]{Lemma}
\newtheorem{proposition}[theorem]{Proposition}
\newtheorem{corollary}[theorem]{Corollary}
\newcommand{\PROVED}{\textcolor{green!45!black}{\textbf{PROVED}}}
\newcommand{\OPEN}{\textcolor{red!70!black}{\textbf{OPEN}}}
\newcommand{\AUDIT}{\textcolor{orange!75!black}{\textbf{AUDIT}}}
\newcommand{\Ran}{\operatorname{Ran}}
\newcommand{\Ker}{\operatorname{Ker}}
\newcommand{\old}{\mathrm{old}}
\newcommand{\Hod}{\mathrm{Hod}}
\newcommand{\comp}{\mathrm{comp}}
\newcommand{\cC}{\mathscr C}
\newcommand{\cZ}{\mathscr Z}
\newcommand{\cL}{\mathscr L}
\title{Row (d): Gamma--Fourier--Tate amplitude audit\\
v90 (critical Euler--Gamma boundary construction)}
\author{}
\date{12 August 2026}

\begin{document}
\maketitle

\begin{abstract}
This note attacks the amplitude identity left by v83--v86 and reconciles the
result with the supplied files \texttt{main.tex} and
\texttt{row-d-local-analysis.tex}.  It constructs a
source-defined doubled logarithmic Tate port and proves that its divergence is
exactly $-B_Q$, thereby closing G39 at every finite threshold.  It also derives
the connection term lost when the Fourier--Gamma identity is differentiated
after a moving compression.  The calculation does not prove the required
positive energy estimate G40.  A new exact balanced factorization of the
localized primitive operator is proved below; it supplies explicit meanings
for the previously abstract factors $X_T,Y_T$.  It also shows precisely why
the local Gamma--Tate normalizations do not by themselves imply the missing
unit contraction.  New critical gcd and Archimedean Gamma GNS boundaries,
an exact Gamma ratio--scale formula deriving $m_0$, and the coherent completed
Euler--Gamma metric are constructed.  A domain theorem rules out unweighted
Gamma deconvolution on every nonzero compact-window source.  The critical
form boundary is then constructed and identified with divisor incidence, the
gcd GNS state and M\"obius whitening.  Explicit Julia, beta--Gamma and Euler
Markov colligations derive $g_\Gamma$, $m_0$ and every prime-power coefficient
as metric/defect rates.  Consequently the remaining bridge is one
operator-level Kato/Julia compression identity; neither condition $G$ nor
row (d) is declared closed.
\end{abstract}

\section{Input from v83--v86}

The following identities are taken as the proved input of the preceding
notes.  At the unweighted endpoint let
\[
 G'=R^*\cL R-\frac12(QG+GQ),\qquad Gv=S\delta_1,
 \qquad Q\delta_d=(\log d)\delta_d .
\]
Let $P_{\old}$ be the old-coordinate projection and let $M_0$ be the old
block of $G$.  Differentiation of the harmonic normal equation gives
\begin{equation}
 P_{\old}G'v=M_0h'.                                      \tag{1}
\end{equation}
The natural Hodge current and its old analysis operator are
\[
 W_{\Hod}=\cL^{1/2}Rv,
 \qquad H_0=\cL^{1/2}RP_{\old}.
\]
Consequently
\begin{equation}
 H_0^*W_{\Hod}=M_0h'+B_Q,
 \qquad B_Q:=\frac12P_{\old}GQv.                         \tag{2}
\end{equation}
In the Zeta gauge, with $m=\cZ^{-1}\delta_1$, this is
\begin{equation}
 B_Q=\frac S2P_{\old}
 \left[-\cC\delta_1+\cZ\cC^*m\right].                  \tag{3}
\end{equation}
The first summand is the directed von--Mangoldt forcing; the second is its
adjoint return through the Zeta triangular system.

The desired archimedean cancellation is
\begin{equation}
 H_{\infty,0}^*W_\infty=-B_Q.                            \tag{6}
\end{equation}

\paragraph{Status.} Equations (1)--(3) are \PROVED{} within the stated
finite-dimensional conventions of v83--v86.  Equation (6) is proved below in
Theorem G39, equations (27e)--(27h).

\section{Exact solvability of the amplitude equation}

The next statement is elementary, but it is the correct non-circular test for
any proposed Gamma--Tate construction.

\begin{theorem}[Amplitude range theorem --- \PROVED]
Let $H:\mathcal K_0\to\mathcal K_\infty$ and
$B:E\to\mathcal K_0$ be bounded operators (all spaces may be finite
dimensional).  The equation
\[
 H^*W=-B,\qquad W:E\to\mathcal K_\infty,
\]
has a solution if and only if
\begin{equation}
 \Ran B\subseteq\Ran H^*.                                \tag{4}
\end{equation}
Equivalently, in finite dimension,
\begin{equation}
 \Ker H\subseteq\Ker B^*.                               \tag{5}
\end{equation}
When it exists, the unique solution whose range is orthogonal to $\Ker H^*$ is
\begin{equation}
 W_{\min}=-(H^*)^\dagger B,                              \tag{7}
\end{equation}
and every solution is $W_{\min}+K$ with $\Ran K\subseteq\Ker H^*$.  Moreover
\begin{equation}
 W^*W\ge W_{\min}^*W_{\min}
 =B^*(H^*H)^\dagger B.                                  \tag{8}
\end{equation}
\end{theorem}

\begin{proof}
The equation is solvable precisely when every column of $B$ lies in
$\Ran H^*$.  The orthogonal-complement identity
$(\Ran H^*)^\perp=\Ker H$ gives (5).  Moore--Penrose inversion gives (7).
The general solution differs from it by a map into $\Ker H^*$; this summand is
orthogonal to $W_{\min}$, proving (8).
\end{proof}

Applied to (6), the first genuinely new target is therefore
\begin{equation}
 \boxed{\Ran B_Q\subseteq\Ran H_{\infty,0}^*.}           \tag{9}
\end{equation}
This is an amplitude/range statement.  It is not implied by equality of the
associated quadratic-form generators.

\section{Why the known Fourier--Gamma identity does not yet give (6)}

The proved Fourier--Gamma bridge has the form
\begin{equation}
 G_\Gamma-m_0I=-(K+X).                                  \tag{10}
\end{equation}
It identifies an archimedean generator.  The desired equation (6), however,
contains an old/new matrix coefficient of a factorization of that generator.
Factorizations are not determined by their Gram operator: if
$G_\Gamma-m_0I=J^*\Sigma J$, then $UJ$ gives the same Gram identity for every
unitary $U$ commuting with the signature, while its compressed amplitudes can
change.

There is a second, independent issue: the relevant harmonic/old projection
moves under the deformation.

\begin{lemma}[Moving-compression connection --- \PROVED]
Let $U(t)$ be a differentiable unitary family, $A=U(0)^*U'(0)$, and
$P(t)=U(t)P U(t)^*$.  Then
\begin{equation}
 P'(0)=U(0)[A,P]U(0)^*.                                  \tag{11}
\end{equation}
For a differentiable operator family $L(t)$,
\begin{align}
 \frac d{dt}\bigg|_0 P(t)L(t)P(t)
  ={}&P'(0)L(0)P(0)+P(0)L'(0)P(0)\notag\\
    &+P(0)L(0)P'(0).                                    \tag{12}
\end{align}
Thus a formula for $L'(0)$ alone does not determine the differentiated
compressed block.
\end{lemma}

\begin{proof}
Differentiate $U(t)PU(t)^*$ and then apply the product rule to the compression.
\end{proof}

In the present problem the two terms containing $P'(0)$ are precisely the
kind of connection terms capable of producing the adjoint-return summand
$\cZ\cC^*m$ in (3).  Therefore a valid derivation of (6) from Tate must specify
the boundary embeddings and prove the differentiated compressed intertwining,
not only cite (10).

\section{A sharp conditional completion}

Assume a source-defined Gamma--Tate factor $H_{\infty,0}$ is given.  Then (9)
constructs, without choice,
\[
 W_{\infty,\min}=-(H_{\infty,0}^*)^\dagger B_Q .
\]
Combining it with (2) yields
\[
 (H_0^*\oplus H_{\infty,0}^*)
 (W_{\Hod}\oplus W_{\infty,\min})=M_0h'.               \tag{13}
\]
This proves the required divergence identity.  To obtain the Pythagorean
storage identity one still needs the independently derived energy budget
\begin{equation}
 S_M'\ge B_Q^*(H_{\infty,0}^*H_{\infty,0})^\dagger B_Q
       +\text{the finite-part minimum energy}.           \tag{14}
\end{equation}
Defining a residual as the square root of the difference in (14) would be
circular unless positivity of that difference has already been proved from
the source network.

\section{Source audit: Hilbertian transport is not yet unitary amplitude}

The semilocal Sonin theorem of Connes--Consani--Moscovici supplies a
hilbertian isomorphism
\[
 \theta_S:{\bf S}_\lambda(\mathbb R,e_\infty)
 \longrightarrow {\bf S}_\lambda(X_S,\alpha).
\]
This is enough to transport the underlying topological Hilbert space and to
prove stability when places are added.  It is not, as stated, an isometry for
the fixed physical metrics entering row (d).  In the same discussion the
authors explicitly note that the finite set $S$ plays a key role in fixing the
inner product of the corresponding space of entire functions.

This distinction has an exact operator consequence.  Put
\begin{equation}
 D_S:=\theta_S^*\theta_S>0.                               \tag{15}
\end{equation}
Then the canonically normalized isometry is
\begin{equation}
 U_S:=\theta_SD_S^{-1/2},\qquad U_S^*U_S=I.              \tag{16}
\end{equation}
For a differentiable deformation $S=S(t)$ one has
\begin{equation}
 U_S'=\theta_S'D_S^{-1/2}
       +\theta_S(D_S^{-1/2})'.                           \tag{17}
\end{equation}
The second summand is a metric-connection term.  It vanishes only if
$D_S'=0$, i.e. only if the raw hilbertian transport is already isometric along
the deformation.

\begin{proposition}[Metric correction identity --- \PROVED]
Let $D(t)=\theta(t)^*\theta(t)>0$ and
$U(t)=\theta(t)D(t)^{-1/2}$.  Then $U(t)$ is isometric and
\begin{equation}
 U^*U'=D^{-1/2}\theta^*\theta'D^{-1/2}
       +D^{1/2}(D^{-1/2})',                              \tag{18}
\end{equation}
whose Hermitian part is zero.  Hence the Hermitian part of the raw logarithmic
connection $D^{-1/2}\theta^*\theta'D^{-1/2}$ is cancelled exactly by the
derivative of the metric normalization.
\end{proposition}

\begin{proof}
Equation (18) follows by the product rule.  Differentiating $U^*U=I$ gives
$U'^*U+U^*U'=0$.
\end{proof}

The correction in (17) is computable without assuming that $D$ and $D'$ commute.
If $X=(D^{-1/2})'$, differentiation of
$D^{-1/2}DD^{-1/2}=I$ gives the Sylvester equation
\begin{equation}
 XD^{1/2}+D^{1/2}X=-D^{-1/2}D'D^{-1/2}.                  \tag{19}
\end{equation}
Since $D^{1/2}>0$, (19) has a unique bounded selfadjoint solution for
selfadjoint $D'$.  Thus the next calculation has no unitary gauge freedom at
the metric-correction level: once the semilocal Gram derivative $D'$ is known,
the missing connection $X$ is uniquely fixed.

Thus the published stability theorem cannot be used to infer a norm-one
storage law without explicitly carrying $D_S^{-1/2}$ and its derivative.
The missing correction is of exactly the same structural type as the moving
harmonic-lift term found in v83--v86.  This rules out the shortcut
``Sonin stability $\Rightarrow$ (6)'', while leaving open the possibility that
the fully normalized connection in (17) produces $-B_Q$.

\section{Exact CCM multiplier and its normalized connection}

The source definitions permit the metric correction to be evaluated
explicitly.  On the multiplicative Fourier line put
\begin{equation}
 a_{S,t}(s):=\prod_{p\in S_f}
 \left(1-p^{-t/2-is}\right).                             \tag{20}
\end{equation}
The dual Sonin map $\theta_{S,t}$ is multiplication by $a_{S,t}$, while the
Zeta map $\eta_{S,t}$ is multiplication by
$\overline{a_{S,t}}^{-1}$ (for real $s$).  Hence
\begin{equation}
 \theta_{S,t}^*\eta_{S,t}=I,
 \qquad
 D_{S,t}=\theta_{S,t}^*\theta_{S,t}=M_{|a_{S,t}|^2}.     \tag{21}
\end{equation}
In particular, the polar normalization is not abstract:
\begin{equation}
 U_{S,t}=\theta_{S,t}D_{S,t}^{-1/2}
 =M_{a_{S,t}/|a_{S,t}|}.                                \tag{22}
\end{equation}

Define the finite Euler logarithmic multiplier
\begin{equation}
 c_{S,t}(s):=\sum_{p\in S_f}(\log p)
 \frac{p^{-t/2-is}}{1-p^{-t/2-is}}.                     \tag{23}
\end{equation}

\begin{theorem}[Normalized Euler connection --- \PROVED]
On any compact $t$-interval on which the factors in (20) do not vanish on the
real Fourier line,
\begin{align}
 \theta_{S,t}^{-1}\partial_t\theta_{S,t}
   &=\frac12M_{c_{S,t}},                                 \tag{24}\\
 D_{S,t}^{-1}D_{S,t}'
   &=\frac12M_{c_{S,t}+\overline{c_{S,t}}},              \tag{25}\\
 U_{S,t}^*U_{S,t}'
   &=\frac14M_{c_{S,t}-\overline{c_{S,t}}}.              \tag{26}
\end{align}
Thus metric normalization cancels exactly the Hermitian Euler component and
retains exactly its skew-Hermitian component.
\end{theorem}

\begin{proof}
For one factor, with $r_p(t)=p^{-t/2}$,
\[
 \partial_t\log(1-r_p(t)p^{-is})
 =\frac{\log p}{2}\frac{p^{-t/2-is}}
 {1-p^{-t/2-is}}.
\]
Summation gives (24).  Equations (25) and (26) follow by differentiating
$|a|^2$ and $a/|a|$, respectively.
\end{proof}

This theorem sharpens the target.  The normalized CCM transport contributes
only the oriented difference $c-\bar c$.  Therefore the symmetric Hodge
energy and the amplitude $B_Q$ cannot be read from this phase connection
alone.  They can arise only after the old/new (Sonin/cutoff) compression and
the two moving boundary embeddings are differentiated.  Equivalently, the
remaining source identity is now the explicit compressed formula
\begin{equation}
 P_{\old}\left(
 j_{\old}'{}^*U^*j_{\new}
 +j_{\old}^*U^*U'j_{\new}
 +j_{\old}^*U^*j_{\new}'
 \right)
 =-B_Q,                                                  \tag{27}
\end{equation}
with $U^*U'$ fixed by (26).  There is no longer an unspecified bulk
connection in (27); the only unknown data are the differentiated Tate/cutoff
boundary maps.

\paragraph{Source boundary.}
The cited CCM theorem proves (at fixed finite set of places $S$ and fixed
cutoff $\lambda$) the stability of the Sonin space and the duality
$\theta_S^*\eta_S=I$.  It does not state a differentiable family of Tate
boundary embeddings $j_{\old}(t),j_{\new}(t)$, nor formula (27).  Consequently
those maps cannot be silently imported from that theorem.  The finite
divisibility model below constructs the required logarithmic boundary port
directly, without attributing it to the fixed-parameter CCM theorem.

\section{Closure of G39: the doubled logarithmic Tate port}

We now give a source-defined realization of the boundary amplitude.  Let
$\mathcal H_N=\operatorname{span}\{\delta_d:1\le d\le N\}$ and let
$S_n\delta_d=\delta_{nd}$ when $nd\le N$ and zero otherwise.  Put
\begin{align}
 \cZ_t&:=\sum_{n\le N}n^{-t/2}S_n,\notag\\
 \cC_t&:=\sum_{n\le N}\Lambda(n)n^{-t/2}S_n,
 \qquad Q\delta_d=(\log d)\delta_d.                     \tag{27a}
\end{align}
All sums and products are in the finite divisibility algebra, so there are no
domain or convergence issues.

\begin{lemma}[Finite logarithmic connection --- \PROVED]
One has
\begin{align}
 [Q,\cZ_t]&=\cZ_t\cC_t=\cC_t\cZ_t,                    \tag{27b}\\
 \cZ_t'&=-\frac12\cZ_t\cC_t,                           \tag{27c}\\
 (\cZ_t^{-1})'&=\frac12\cC_t\cZ_t^{-1}.                \tag{27d}
\end{align}
Thus $\cC_t=-2\cZ_t^{-1}\cZ_t'$ is the logarithmic
derivative of the Zeta transport, and $\cC_t^*$ is its adjoint-return
connection.
\end{lemma}

\begin{proof}
Since $[Q,S_n]=(\log n)S_n$, the coefficient of $S_k$ in
$[Q,\cZ_t]$ is $(\log k)k^{-t/2}$.  The coefficient of $S_k$ in
$\cZ_t\cC_t$ is
\[
 k^{-t/2}\sum_{d\mid k}\Lambda(d)
 =k^{-t/2}\log k.
\]
This proves (27b).  Direct differentiation gives (27c), and differentiating
$\cZ_t^{-1}\cZ_t=I$ gives (27d).
\end{proof}

At the critical endpoint suppress the subscript $t$, put
\[
 m=\cZ^{-1}\delta_1,\qquad
 \mathcal K_\infty^{(1)}:=\mathcal H_N\oplus\mathcal H_N,
\]
and define, before invoking any sign condition,
\begin{align}
 W_\infty
 &:=\sqrt{\frac S2}
 \binom{\cC\delta_1}{\cC^*m},                           \tag{27e}\\
 H_{\infty,0}x
 &:=\sqrt{\frac S2}
 \binom{P_{\old}^*x}{-\cZ^*P_{\old}^*x}.              \tag{27f}
\end{align}
The first component of (27e) is the differentiated Zeta/Möbius incoming
port; the second is its adjoint return.  Formula (27f) is the graph analysis
of the dual Tate termination $[I,-\cZ]$.

\begin{theorem}[G39, exact Gamma--Tate amplitude --- \PROVED]
The doubled logarithmic port (27e)--(27f) is source-defined by the finite
Euler transport and satisfies
\begin{equation}
 \boxed{H_{\infty,0}^*W_\infty=-B_Q}.                   \tag{27g}
\end{equation}
Equivalently, it supplies exactly the missing divergence in (2), and hence
\begin{equation}
 (H_0^*\oplus H_{\infty,0}^*)
 (W_{\Hod}\oplus W_\infty)=M_0h'.                       \tag{27h}
\end{equation}
No positivity of $G$, no Moore--Penrose choice, and no residual square root is
used in this construction.
\end{theorem}

\begin{proof}
Taking the adjoint of (27f) and applying it to (27e) gives
\[
 H_{\infty,0}^*W_\infty
 =\frac S2P_{\old}
 \left(\cC\delta_1-\cZ\cC^*m\right).
\]
By the already proved v84 formula (3), the right-hand side is exactly $-B_Q$.
Adding this equality to (2) proves (27h).
\end{proof}

The construction also explains the two moving-boundary contributions in
(27): they are the two orientations of the logarithmic derivative of the
dual pair $\cZ_t^{-1}$ and $\cZ_t^{-*}$.  The bulk normalized phase accounts
for their skew part, while the graph termination $[I,-\cZ]$ retains their
oriented difference.  Thus G39 is closed at every finite threshold.  This
does not assert the Hilbert-space energy bound G40.

\section{Reverse-engineered completed boundary object}

We now construct the object required by the proof without presupposing its
positivity.  This separates existence from the load-bearing sign theorem.
Let
\[
 A_0=\mathcal G_0^*\mathcal G_0\ge0,
 \qquad q=M_0h',\qquad \beta=S_M'.
\]
All operators below are understood on the finite threshold space; the same
formulas hold for closed ranges in the Hilbert-space setting.

\begin{theorem}[Canonical completed boundary colligation --- \PROVED]
Assume $q\in\Ran A_0^{1/2}$ and define
\begin{equation}
 w_0:=\mathcal G_0A_0^\dagger q,
 \qquad
 \Delta_T:=\beta-q^*A_0^\dagger q.                      \tag{28}
\end{equation}
Let $\Delta_T=J_T|\Delta_T|$ be its polar decomposition, where
$J_T=\operatorname{sgn}(\Delta_T)$ on
$\overline{\Ran|\Delta_T|}$, and put
\begin{equation}
 \mathcal K_T^{\rm bd}:=\overline{\Ran|\Delta_T|},
 \qquad r_T:=|\Delta_T|^{1/2}.                           \tag{29}
\end{equation}
Then the map
\begin{equation}
 \mathcal W_Te:=w_0e\oplus r_Te                         \tag{30}
\end{equation}
into the Krein boundary space
$\overline{\Ran\mathcal G_0}\oplus(\mathcal K_T^{\rm bd},J_T)$
satisfies the two exact identities
\begin{align}
 \mathcal G_0^*P_0\mathcal W_T&=q,                      \tag{31}\\
 \mathcal W_T^*(I\oplus J_T)\mathcal W_T&=\beta.        \tag{32}
\end{align}
It is minimal: any other realization of (31)--(32) contains this one up to a
Krein-isometric summand orthogonal to the generated boundary space.
\end{theorem}

\begin{proof}
Since $q\in\Ran A_0^{1/2}$, Moore--Penrose calculus gives
\[
 \mathcal G_0^*w_0
 =A_0A_0^\dagger q=q,
 \qquad
 w_0^*w_0=q^*A_0^\dagger q.
\]
Moreover $r_T^*J_Tr_T=\Delta_T$.  Adding the last two identities proves
(31)--(32).  Minimality follows from the minimum-norm property of $w_0$ and
the minimal polar realization of the selfadjoint form $\Delta_T$.
\end{proof}

This is the desired inverse-engineered object: it always exists once the
previously proved range condition holds, but in general its boundary metric is
Krein rather than Hilbert.

\begin{theorem}[Exact positivity threshold --- \PROVED]
The following assertions are equivalent:
\begin{enumerate}
\item the canonical boundary metric in (29) is positive, $J_T=I$;
\item $\Delta_T\ge0$;
\item there exist Hilbert-space operators $W,R$ such that
\[
 \mathcal G_0^*W=q,
 \qquad W^*W+R^*R=\beta;
\]
\item the block completion is positive,
\begin{equation}
 \begin{pmatrix}A_0&q\\q^*&\beta\end{pmatrix}\ge0;      \tag{33}
\end{equation}
\item the v85 Pythagorean identity, hence the corresponding threshold case of
condition $G$, holds.
\end{enumerate}
\end{theorem}

\begin{proof}
The equivalence of (1) and (2) is the definition of $J_T$.  If (2) holds take
$W=w_0$ and $R=\Delta_T^{1/2}$.  Conversely minimum norm gives
$W^*W\ge q^*A_0^\dagger q$, so any Hilbert realization implies
$\Delta_T\ge0$.  The equivalence with (33) is the generalized Schur-complement
criterion, including $q\in\Ran A_0^{1/2}$.  The last equivalence is precisely
the reduction proved in v85--v86.
\end{proof}

The theorem shows both the power and the limit of pure inverse engineering.
It constructs the unique missing port and its exact metric, but proving that
the metric is Hilbert rather than Krein is exactly the remaining content of
$G$.  Renaming $r_T$ a ``Tate residual'' would therefore be circular.

\subsection{The genuinely new theorem sufficient to close row (d)}

The preceding construction identifies a source-level statement which is
strictly stronger than a formal completion and is independently falsifiable:

\begin{quote}
\textbf{Tate boundary positivity theorem (OPEN).}
There is an isomorphism
\[
 \Phi_T:\mathcal K_T^{\rm bd}\longrightarrow
 \mathcal T_T^{\rm out}
\]
onto a boundary space obtained from the two terminated Tate evaluations and
the Gamma storage tail, constructed before using $\Delta_T$, such that
\begin{equation}
 \langle\Phi_Tx,\Phi_Ty\rangle_{\mathcal T_T^{\rm out}}
 =\langle J_Tx,y\rangle.                                \tag{34}
\end{equation}
The target carries its ordinary positive $L^2$/Gamma-storage inner product.
\end{quote}

If (34) is proved from Fourier--Gamma--Tate, then its left-hand side is
positive, hence $J_T=I$ on the minimal boundary space.  The preceding theorem
then gives $G$ immediately.  Unlike defining
$R=\Delta_T^{1/2}$, statement (34) is non-circular only when
$\mathcal T_T^{\rm out}$ and $\Phi_T$ are given by explicit source formulas
which do not use $\Delta_T$ or its sign.

The proved amplitude identity (27g) is the divergence part of (34).  The missing
norm identity is
\begin{equation}
 \Phi_T^*\Phi_T=J_T                                    \tag{35}
\end{equation}
on the generated boundary space.  Thus (27g) alone is necessary but not
sufficient; (35) is exactly what prevents an indefinite completion from being
mistaken for a conservative Hilbert-space network.

\section{Intermediate status ledger (superseded by the final ledger below)}

\subsection*{PROVED}
\begin{itemize}
\item The v83 Möbius/von--Mangoldt forcing and the v84 excess divergence
      formula (2)--(3), subject to their stated finite Zeta-gauge definitions.
\item The range/kernel criterion (4)--(5), the canonical current (7), and its
      exact minimum energy (8).
\item The moving-compression connection formula (11)--(12).
\item The metric-normalization correction (15)--(18).
\item The exact CCM multiplier identities (20)--(26).
\item G39: the source-defined doubled logarithmic port (27e)--(27f), the exact
      amplitude identity (27g), and the completed divergence (27h).
\item The canonical Krein boundary colligation (28)--(32) and the exact
      equivalence between its positivity, block positivity (33), the v85
      storage identity, and $G$.
\item The amplitude range inclusion (9) for the constructed port, since
      $-B_Q=H_{\infty,0}^*W_\infty$, and the composite divergence (13)/(27h).
\end{itemize}

\subsection*{OPEN}
\begin{itemize}
\item The non-circular source energy balance (14), hence the v85 Pythagorean
      identity, condition $G$, and row (d).
\item The Tate boundary positivity theorem (34)--(35), with a source-defined
      positive target and an intertwiner independent of $\Delta_T$.
\end{itemize}

\subsection*{AUDIT / forbidden shortcuts}
\begin{itemize}
\item Equation (10) is an operator-generator identity, not an amplitude
      identity; it does not by itself imply (6).
\item Möbius/Zeta adjoint-inverse triangularity does not determine the moving
      old/new connection block.
\item A hilbertian Sonin isomorphism is not automatically an isometry in the
      row-(d) physical metric.  The derivative of $D_S^{-1/2}$ cannot be
      discarded.
\item The parameter deformation and the differentiated Tate boundary maps are
      not supplied by the fixed-$S$, fixed-$\lambda$ Sonin stability theorem.
\item Choosing $W_\infty$ by Moore--Penrose inversion is legitimate only after
      (9); claiming the desired budget for that choice is equivalent to the
      remaining inequality and cannot serve as its proof.
\item The inverse-engineered boundary space exists canonically, but its metric
      is a priori Krein.  Declaring it positive or calling
      $|\Delta_T|^{1/2}$ a Hilbert residual would assume $G$.
\item No statement in this note declares $G$ or row (d) proved.
\end{itemize}

\section{Next exact calculation}
With G39 closed, the next calculation is exclusively the G40 energy theorem.
For the explicitly constructed port it asks whether its source norm, together
with the finite Hodge current and the Gamma storage tail, telescopes to
\[
 S_M'=(M_0h')^*A_0^\dagger(M_0h')+R_T^*R_T
\]
with $R_T$ constructed before assuming the sign of the difference.  The
divergence is no longer an unknown.  The remaining task is to identify the
positive Gamma--Tate storage realization (34)--(35); this is G40, not G39.

\section{G40 bridge audit: KYP completion and the circularity barrier}

The completion of G39 supplies an exact current, but amplitude and passivity
are logically different.  The following theorem determines, without guesswork,
the least energy compatible with the proved amplitude.

Put $H:=H_0\oplus H_{\infty,0}$ and
$q:=M_0h'$.  By (27h), $q\in\Ran H^*$.  Define
\begin{equation}
 W_{\min}:=(H^*)^\dagger q,\qquad
 E_{\min}:=q^*(H^*H)^\dagger q.                         \tag{36}
\end{equation}

\begin{theorem}[Minimal passive completion test --- \PROVED]
For the source budget $\beta=S_M'$, the following are equivalent:
\begin{enumerate}
\item there are a Hilbert-space current $W$ and residual $R$ satisfying
\[
 H^*W=q,\qquad \beta=W^*W+R^*R;
\]
\item $\beta-E_{\min}\ge0$;
\item the KYP block
\begin{equation}
 \mathfrak K_T:=
 \begin{pmatrix}H^*H&q\\q^*&\beta\end{pmatrix}          \tag{37}
\end{equation}
is positive.
\end{enumerate}
If these conditions hold, the canonical solution is
\begin{equation}
 W=W_{\min},\qquad R=(\beta-E_{\min})^{1/2}.             \tag{38}
\end{equation}
\end{theorem}

\begin{proof}
Moore--Penrose minimality gives $W^*W\ge E_{\min}$ for every solution of
$H^*W=q$.  Hence a passive completion implies (2).  Conversely (2) gives
(38).  Equivalence with (3) is the generalized Schur-complement theorem,
using $q\in\Ran H^*$, which is already supplied by G39.
\end{proof}

This theorem is a useful advance: G39 has discharged the entire range part of
the generalized Schur criterion.  G40 consists only of the remaining scalar
operator inequality
\begin{equation}
 \boxed{S_M'-q^*(H^*H)^\dagger q\ge0.}                  \tag{39}
\end{equation}
However, inverse engineering cannot prove (39) by defining the residual in
(38): the square root exists as a Hilbert operator if and only if (39) is
already true.

\begin{theorem}[No-free-bridge theorem --- \PROVED]
Suppose a proposed G40 bridge is constructed only from the already fixed data
$(H,q,\beta)$ and is required to satisfy the exact amplitude and energy
identities.  Then its existence as a Hilbert-space colligation is equivalent
to (39).  In particular, polar decomposition, Julia completion, graph
normalization, Moore--Penrose lifting, or adjoining a formal residual cannot
establish G40 unless one independently proves positivity of (39).
\end{theorem}

\begin{proof}
Every such bridge gives a pair $(W,R)$ in item (1) of the preceding theorem,
so it implies (39).  Conversely (39) constructs the bridge by (38).  Each of
the listed procedures is a realization of this same completion and therefore
cannot change the equivalence.
\end{proof}

Thus the non-circular content still required from Gamma--Tate is an
\emph{independently positive} source formula for the left side of (39), for
example
\begin{equation}
 S_M'-q^*(H^*H)^\dagger q
 =\int_{\Omega_T}K_T(\omega)^*K_T(\omega)\,d\mu_T(\omega)
  +E_{0,T}^*E_{0,T}+E_{1,T}^*E_{1,T},                  \tag{40}
\end{equation}
where $K_T,E_{0,T},E_{1,T}$ must be given explicitly from the Gamma tail and
the two Tate evaluations without using the left side or its spectral
decomposition.  Formula (40), not a formal completion, is the precise new
mathematical theorem which would prove G40.

\subsection{Compatibility with the fixed old factor}

There is one indispensable bookkeeping point.  The analysis operator $H$ in
(36) is the explicit amplitude realization of G39.  The old factor used in
the threshold induction is fixed independently by
\[
 A_0=\mathcal G_0^*\mathcal G_0.
\]
One may replace $\mathcal G_0$ by $H$ in the storage argument only after
proving the Gram compatibility
\begin{equation}
H^*H=A_0,                                               \tag{41}
\end{equation}
or, more generally, constructing a source-defined isometry between their
minimal factorizations.  G39 proves $H^*W=q$; it does not by itself prove
(41).

The notation must be reconciled with the supplied row (d).  There
$R_0=X_0^*X_0$ is the \emph{positive reference} and
$L_0=Y_0^*Y_0$ is the opposite channel; neither $X_0$ alone nor $Y_0$ alone
factors the signed old block.  Under the old-core hypothesis set
\begin{equation}
 T_0=R_0^{\dagger/2}L_0R_0^{\dagger/2},
 \qquad D_0=I-T_0.                                    \tag{41a}
\end{equation}
On the support of $R_0$, one has the exact form identity
\begin{equation}
 A_0=R_0-L_0
 =R_0^{1/2}D_0R_0^{1/2}.                              \tag{41b}
\end{equation}
Consequently, if and only if $D_0\ge0$, a minimal source factor of the signed
old block is
\begin{equation}
 \mathcal G_0=D_0^{1/2}R_0^{1/2},
 \qquad \mathcal G_0^*\mathcal G_0=A_0,               \tag{41c}
\end{equation}
up to the standard support projections.  Thus the actual compatibility test
for the G39 current is (41) with $\mathcal G_0$ from (41c), not
$H^*H=X_0^*X_0$.

\begin{proof}
Since $L_0=R_0^{1/2}T_0R_0^{1/2}$ on
$\overline{\Ran R_0}$, subtraction gives (41b).  The old-core hypothesis is
$0\le T_0\le I$, hence $D_0\ge0$ and functional calculus gives (41c).
Conversely existence of a Hilbert factor of (41b) forces the signed old block
to be nonnegative.
\end{proof}

This correction also locates the return-load condition in the current
language.  A normalized old load can be lifted through $\mathcal G_0^*$
exactly when it lies in $\Ran D_0^{1/2}$ after reference normalization.  This
is the range condition explicitly left open before (and used conditionally
in) \texttt{eq:returnsum} of the supplied row (d).  Hence replacing
$\mathcal G_0$ by the reference factor $X_0$ would erase precisely the
capacity defect that G40 is required to control.

\begin{theorem}[Exact reference-harmonic capacity identity --- \PROVED]
Write the old/new columns of the authoritative factors as
\[
 X=(X_0,X_E),\qquad Y=(Y_0,Y_E),
\]
and put
\[
 R_0=X_0^*X_0,\quad L_0=Y_0^*Y_0,\quad
 r=X_0^*X_E,\quad \ell=Y_0^*Y_E.
\]
Let $H=R_0^\dagger r$ and assume the usual reference range condition
$r\in\Ran R_0$.  Define
\begin{align}
 S_E&:=X_E^*X_E-r^*R_0^\dagger r,                     \tag{41d}\\
 Z_E&:=Y_E-Y_0H,                                      \tag{41e}\\
 b_E&:=L_0H-\ell.                                     \tag{41f}
\end{align}
Then the reference-harmonic column $v_E=(-H,I)^t$ satisfies
\begin{align}
 v_E^*(X^*X-Y^*Y)v_E&=S_E-Z_E^*Z_E,                  \tag{41g}\\
 P_{\rm old}(X^*X-Y^*Y)v_E&=b_E.                     \tag{41h}
\end{align}
If $A_0=R_0-L_0\ge0$, the full old/new signed block is nonnegative if and
only if
\begin{equation}
 b_E\in\Ran A_0^{1/2},\qquad
 \boxed{S_E-Z_E^*Z_E-b_E^*A_0^\dagger b_E\ge0.}       \tag{41i}
\end{equation}
\end{theorem}

\begin{proof}
The normal equation $R_0H=r$ gives
$X_0^*(X_E-X_0H)=0$ and
$(X_E-X_0H)^*(X_E-X_0H)=S_E$.  This proves the reference part of
(41g); its opposite part is $Z_E^*Z_E$ by definition.  The old cross of the
signed block against $v_E$ is
\[
 (r-\ell)-(R_0-L_0)H=L_0H-\ell=b_E,
\]
which proves (41h).  Congruence by the invertible triangular change from the
coordinate newborn column to $v_E$ leaves the old block equal to $A_0$ and
changes the other entries to (41h) and (41g).  The generalized Schur
criterion gives exactly (41i).
\end{proof}

Formula (41i) is the authoritative version of the G40 storage target.  It
also proves why the derivative calculation in v83 found a moving-lift square:
$v_E$ is harmonic for $R_0$, but its signed old divergence is the generally
nonzero $b_E$.  The Moore--Penrose term in (41i) is the energy required to
remove that divergence.  Any proposed Hodge residual must be identified with
the entire left side of (41i), not merely with $S_E-Z_E^*Z_E$.

On supported vectors for which $S_E^{-1/2}$ is defined, the normalized
variables of the supplied row (d) obey the exact translation
\begin{equation}
 u_e=D_0h_e-q_e
 =-R_0^{\dagger/2}b_ES_E^{\dagger/2}e.                \tag{41j}
\end{equation}
Indeed, substitute $R_0H=r$ into the definitions of $D_0,h_e,q_e$.  Thus the
apparently different range obstruction in \texttt{eq:returnsum} is precisely
the normalized form of the Feshbach range condition in (41i), and its return
energy is the normalized Moore--Penrose term there.  This closes the
bookkeeping equivalence between the threshold-capacity and balanced-factor
languages; it does not prove their common sign.

\begin{theorem}[Physical old-compatible amplitude port --- \PROVED]
Assume the inductive old-core hypothesis $A_0\ge0$ and the range condition
$b_E\in\Ran A_0^{1/2}$.  With $\mathcal G_0$ as in (41c), define
\begin{equation}
 W_{0,E}:=\mathcal G_0A_0^\dagger b_E.                 \tag{41k}
\end{equation}
Then
\begin{align}
 \mathcal G_0^*W_{0,E}&=b_E,                           \tag{41l}\\
 W_{0,E}^*W_{0,E}&=b_E^*A_0^\dagger b_E.              \tag{41m}
\end{align}
Moreover (41m) is the least Hilbert-current energy among all solutions of
$\mathcal G_0^*W=b_E$.
\end{theorem}

\begin{proof}
Moore--Penrose calculus and the range hypothesis give
$A_0A_0^\dagger b_E=b_E$, proving (41l).  Using
$\mathcal G_0^*\mathcal G_0=A_0$ gives (41m).  Every other solution differs
from (41k) by a vector in $\ker\mathcal G_0^*$, orthogonal to the minimal
solution, proving minimality.
\end{proof}

Thus amplitude and old-factor compatibility are completely realized inside
the authoritative threshold model once the inductive old sign and its range
condition are available.  The remaining newborn energy after paying this
minimal old current is exactly the boxed form (41i).  Identifying the finite
divisibility G39 port with (41k) would require a source-defined unitary (or
isometry plus null channels) preserving both the old Gram and the cross; the
present files provide no such comparison map.

\begin{theorem}[Joint-Gram identification criterion --- \PROVED]
Let $(K,W):E_0\oplus E\to\mathcal K$ and
$(\widetilde K,\widetilde W):E_0\oplus E\to\widetilde{\mathcal K}$ be two
old/new analysis pairs.  There is a unique unitary
\begin{equation}
 U:\overline{\Ran(K,W)}\longrightarrow
   \overline{\Ran(\widetilde K,\widetilde W)},
 \qquad UK=\widetilde K,\quad UW=\widetilde W,          \tag{41n}
\end{equation}
if and only if their complete joint Grams agree:
\begin{equation}
 \begin{pmatrix}K^*K&K^*W\\W^*K&W^*W\end{pmatrix}
 =
 \begin{pmatrix}
 \widetilde K^*\widetilde K&\widetilde K^*\widetilde W\\
 \widetilde W^*\widetilde K&\widetilde W^*\widetilde W
 \end{pmatrix}.                                       \tag{41o}
\end{equation}
\end{theorem}

\begin{proof}
Necessity follows from $U^*U=I$.  Conversely define
$U(Kx+We)=\widetilde Kx+\widetilde We$.  Equality (41o) shows that this rule
is well defined and preserves every inner product.  It extends to an
isometry of the closed generated ranges and is onto by construction.
\end{proof}

For G39 versus the physical threshold port one must first construct source
maps $\iota_0:E_0^{\rm phys}\to E_0^{\rm div}$ and
$\iota_E:E^{\rm phys}\to E^{\rm div}$.  After these have been specified,
(41o) requires simultaneously
\begin{align}
 \iota_0^*H_{\infty,0}^*H_{\infty,0}\iota_0&=A_0,   \tag{41p}\\
 \iota_0^*(-B_Q)\iota_E&=b_E,                         \tag{41q}\\
 \iota_E^*W_\infty^*W_\infty\iota_E
 &=b_E^*A_0^\dagger b_E.                               \tag{41r}
\end{align}
The adjoint cross is automatic once (41q) holds.  Equation (27g) establishes
the left cross in its finite divisibility model, but the source maps and none
of (41p)--(41r) are constructed in the supplied continuum row (d): the two
models use different source spaces and different deformation variables.  In
particular, (42) expands the untransported G39 old Gram, while (41a)--(41c)
expand the physical old Gram; their equality after transport is not a
normalization convention.  This joint-Gram test is
the exact, correctly typed interface required before G39 may be used as the
physical G40 current.

\begin{theorem}[Finite-port/full-Hodge dimension obstruction --- \PROVED]
Suppose the physical old block $A_0$ has infinite rank (in particular, suppose
it is strictly positive on an infinite-dimensional old primitive core).  No
map $\iota_0:E_0^{\rm phys}\to\mathcal H_N$ can satisfy (41p), because
$H_{\infty,0}$ in (27f) has finite-dimensional codomain.  Consequently the
finite divisibility G39 port cannot be unitarily identified with the complete
physical old/Hodge factor of the supplied row (d).
\end{theorem}

\begin{proof}
The rank of
$\iota_0^*H_{\infty,0}^*H_{\infty,0}\iota_0$ is at most
$\dim\mathcal K_\infty^{(1)}=2N$, whereas $A_0$ has infinite rank.  Thus
(41p) is impossible.  Under strict positivity the contradiction can also be
seen from injectivity: (41p) would force $\iota_0$ to be injective into a
finite-dimensional space.
\end{proof}

Hence Theorem G39 remains a valid finite arithmetic amplitude identity, but
it is not a factorization of the full continuum row-(d) current.  It may be
used for compatible finite compressions only.  A full G40 proof requires an
infinite-dimensional continuum port constructed directly from (81)--(84),
not an isometric embedding of (27e)--(27f).

For the doubled logarithmic port its additional old Gram is explicitly
\begin{equation}
 H_{\infty,0}^*H_{\infty,0}
 =\frac S2P_{\old}(I+\cZ\cZ^*)P_{\old}^*.               \tag{42}
\end{equation}
Thus (41) is a concrete, falsifiable operator identity, not a normalization
convention.

The complete G40 bridge can be stated without subdivision as the single
source factorization
\begin{equation}
 \boxed{
 \begin{pmatrix}A_0&q\\q^*&S_M'\end{pmatrix}
 =\mathfrak T_T^*\mathfrak T_T,}                        \tag{43}
\end{equation}
where the old/new columns of $\mathfrak T_T$ must restrict to the G39
amplitude realization and the remaining rows must be explicit Gamma-tail and
Tate-terminal outputs.  Equality (43) simultaneously enforces the correct old
Gram, the proved cross amplitude, and the positive source budget.

Reverse engineering always produces a Krein factorization of the left side of
(43), but it produces a Hilbert factorization if and only if that block is
positive.  By the generalized Schur theorem this is equivalent to
$S_M'-q^*A_0^\dagger q\ge0$, namely G42/G.  Therefore (43) proves G40 only if
$\mathfrak T_T$ is constructed independently from Gamma--Tate formulas; taking
a square root or polar factor of the block would be circular.

\section{The unitary midpoint bridge}

There is a non-circular mechanism capable of proving (43): the midpoint split
of a source-defined unitary Poisson--Fourier operator.  It avoids defining a
Schur residual.

\begin{theorem}[Unitary midpoint storage --- \PROVED]
Let $\mathcal U$ be unitary on a Hilbert space $\mathcal K$, let $W$ be a
current, and define
\begin{equation}
 W_+:=\frac12(I+\mathcal U)W,
 \qquad
 W_-:=\frac12(I-\mathcal U)W.                           \tag{44}
\end{equation}
Then
\begin{equation}
 W^*W=W_+^*W_++W_-^*W_-.                               \tag{45}
\end{equation}
If an analysis operator $H$ satisfies
\begin{equation}
 H^*W=q+B,\qquad H^*W_-=B,                             \tag{46}
\end{equation}
then
\begin{equation}
 H^*W_+=q                                               \tag{47}
\end{equation}
and (45) is the required non-circular positive storage identity.
\end{theorem}

\begin{proof}
Expanding the two squares in (44), the mixed terms cancel and
\[
 W_+^*W_++W_-^*W_-
 =\frac12W^*W+\frac12W^*\mathcal U^*\mathcal UW=W^*W.
\]
Equation (47) follows by subtracting the two equations in (46).
\end{proof}

For the present problem take $W=W_{\Hod}$, so that $W^*W=S_M'$.  G39 has
constructed a current with divergence $-B_Q$.  Consequently the single bridge
identity which would finish G40 is
\begin{equation}
 \boxed{
 W_\infty=-W_-
 =-\frac12(I-\mathcal U_{\rm PF})W_{\Hod},}             \tag{48}
\end{equation}
where $\mathcal U_{\rm PF}$ must be the normalized Poisson--Fourier transport
in the same current Hilbert space.  Indeed, (2) and G39 give
\[
 H^*W_{\Hod}=q+B_Q,\qquad H^*(-W_\infty)=B_Q.
\]
Thus (44)--(48) yield
\begin{align}
 H^*W_+&=q,\notag\\
 S_M'&=W_+^*W_++W_-^*W_-,                              \tag{49}
\end{align}
which is exactly the desired Hilbert storage theorem; $W_-$ is the positive
Tate residual and no square root of $\Delta_T$ has been introduced.

The differentiated Poisson identity from the preceding row-(c) work,
\begin{equation}
 (C+C^\sharp-K-X)f=(\mathcal U_{\rm PF}-I)Kf,           \tag{50}
\end{equation}
has precisely the algebraic form required by (48): the two Euler orientations
give the doubled G39 port and $K+X$ is converted to Gamma by the proved
Fourier--Gamma bridge.  Hence (50) proves (48) at the level of unnormalized
source coefficients.

The remaining referee-level requirement is metric: the map carrying $Kf$ to
$W_{\Hod}$ and the left side of (50) to the doubled current (27e) must be shown
to be isometric for the physical current norm.  The fixed-parameter CCM
theorem gives a hilbertian isomorphism, not this isometry; its normalization
derivative is (25)--(26).  Therefore the exact final identity to check is
\begin{equation}
 \mathcal I_T(\mathcal U_{\rm PF}-I)Kf
 =(\mathcal U_{\rm cur}-I)\mathcal I_TKf,\qquad
 \mathcal I_T^*\mathcal I_T=I,                          \tag{51}
\end{equation}
with $\mathcal I_TKf=W_{\Hod}$.  If (51) is established from the source
normalizations, (48)--(49) prove G40 completely.  Without the second equality
in (51), replacing a hilbertian transport by an isometry would assume the
missing energy statement.

\section{Restricted polar normalization on the Sonin space}

The metric part of (51) can in fact be obtained directly from the available
CCM normalizations, provided the polar normalization is performed after
restriction to the Sonin space.

Let $\mathcal S_T$ be the archimedean Sonin space at the fixed cutoff and
$\mathcal S_T^S$ its semilocal counterpart.  Write
\[
 \Theta_T:=\theta_S|_{\mathcal S_T}:\mathcal S_T
 \longrightarrow\mathcal S_T^S.
\]
The CCM stability theorem says that $\Theta_T$ is bounded and bijective with
bounded inverse, and the Fourier intertwining says
\begin{equation}
 \mathcal F_S\Theta_T=\Theta_T\mathcal F_T.              \tag{52}
\end{equation}
Define the restricted Gram and its polar normalization by
\begin{equation}
 D_T:=\Theta_T^*\Theta_T>0,\qquad
 \mathcal I_T:=\Theta_TD_T^{-1/2}.                       \tag{53}
\end{equation}

\begin{theorem}[Sonin polar isometry --- \PROVED]
The operator $D_T$ commutes with the restricted Fourier symmetry
$\mathcal F_T$, the operator $D_T^{-1/2}$ preserves $\mathcal S_T$, and
\begin{align}
 \mathcal I_T^*\mathcal I_T&=I_{\mathcal S_T},           \tag{54}\\
 \mathcal F_S\mathcal I_T&=\mathcal I_T\mathcal F_T.    \tag{55}
\end{align}
Thus $\mathcal I_T$ is a source-defined unitary from $\mathcal S_T$ onto
$\mathcal S_T^S$; no positivity statement from row (d) is used.
\end{theorem}

\begin{proof}
Both Fourier transforms are unitary and (52) holds.  Hence
\[
 \mathcal F_T^*D_T\mathcal F_T
 =\mathcal F_T^*\Theta_T^*\Theta_T\mathcal F_T
 =\Theta_T^*\mathcal F_S^*\mathcal F_S\Theta_T
 =D_T.
\]
Therefore $D_T$ and every bounded Borel function of $D_T$, in particular
$D_T^{-1/2}$, commute with $\mathcal F_T$.  Since $D_T$ is an operator on the
already restricted space $\mathcal S_T$, its functional calculus preserves
that space.  Equation (54) follows directly:
\[
 \mathcal I_T^*\mathcal I_T
 =D_T^{-1/2}\Theta_T^*\Theta_TD_T^{-1/2}=I.
\]
Finally, using (52) and commutation with $D_T^{-1/2}$,
\[
 \mathcal F_S\mathcal I_T
 =\mathcal F_S\Theta_TD_T^{-1/2}
 =\Theta_T\mathcal F_TD_T^{-1/2}
 =\mathcal I_T\mathcal F_T.
\]
Surjectivity follows from the invertibility of $\Theta_T$ and $D_T^{-1/2}$.
\end{proof}

Consequently the second equality in (51) and the unitary intertwining portion
of the first equality are proved.  The normalization is canonical: it is the
unitary factor in the polar decomposition of the restricted CCM map.

\paragraph{Remaining identification.}
To apply (54)--(55) to (48), one must still prove that the source tangent vector
$Kf$ and the physical Hodge current use these restricted Sonin coordinates:
\begin{equation}
 \boxed{\mathcal I_TKf=W_{\Hod}.}                        \tag{56}
\end{equation}
Equation (56) is not the normalization identity; it is the observation
identity between the Poisson tangent feature and the finite-threshold Hodge
feature.  The multiplier/connection calculations (24)--(27) and G39 prove its
logarithmic boundary component, but equality of the complete feature vectors
must be checked from the definitions of $K$, $R$, and $\cL$.  Thus the
normalization requested in (51) is closed, while the application to the
specific Hodge current retains the explicit observation check (56).

\section{Type audit of the final observation}

As written, (56) suppresses two different codomains.  The polar map has type
\[
 \mathcal I_T:\mathcal S_T\longrightarrow\mathcal S_T^S,
\]
whereas the Hodge feature has type
\[
 W_{\Hod}:E_T\longrightarrow
 \mathcal K_{\Hod,T}:=\overline{\Ran\cL^{1/2}R}.
\]
Thus a referee-level statement requires an additional map
\begin{equation}
 J_T:\mathcal S_T^S\longrightarrow\mathcal K_{\Hod,T}  \tag{57}
\end{equation}
and the correctly typed observation is
\begin{equation}
 J_T\mathcal I_TK_T=W_{\Hod}.                            \tag{58}
\end{equation}

\begin{theorem}[Unitary observation criterion --- \PROVED]
Let $A:E\to\mathcal H_1$ and $B:E\to\mathcal H_2$ be bounded.  There is an
isometry $J:\overline{\Ran A}\to\overline{\Ran B}$ satisfying $JA=B$ if and
only if
\begin{equation}
 A^*A=B^*B.                                              \tag{59}
\end{equation}
In that case $J$ is uniquely determined on $\overline{\Ran A}$ and is unitary
onto $\overline{\Ran B}$.
\end{theorem}

\begin{proof}
Necessity follows from $A^*J^*JA=A^*A$.  Conversely, define $J(Ae)=Be$.
Equation (59) shows first that this is well defined and then that it preserves
inner products (by polarization).  It therefore extends uniquely to an
isometry of the closed ranges, and its range is dense in and closed in
$\overline{\Ran B}$.
\end{proof}

Apply the theorem with $A=\mathcal I_TK_T$ and $B=W_{\Hod}$.  Since
$\mathcal I_T$ is isometric, (59) becomes
\begin{equation}
 \boxed{K_T^*K_T=W_{\Hod}^*W_{\Hod}=S_M'.}              \tag{60}
\end{equation}
Consequently the CCM polar normalization proves the isometry inside the Sonin
transport, but it cannot by itself prove the observation into the Hodge
current space.  The latter exists isometrically precisely when the source
tangent Gram is the physical Hodge Gram.

Equation (60) is a concrete coefficient identity to be proved from the actual
definitions of $K_T$, $R$, and $\cL$:
\begin{equation}
 \boxed{K_T^*K_T=v^*R^*\cL Rv.}                         \tag{61}
\end{equation}
The current local dossier specifies the right side and the logarithmic
boundary of the left side, but it does not contain a definition of $K_T$ as an
operator into the restricted Sonin space sufficient to expand the left side.
Therefore identifying $J_T$ by inverse engineering before proving (61) would
simply define an isometry whose existence is equivalent to the desired Gram
identity.

\section{Audit of the symbol $K$: the midpoint candidate fails}

The preceding row-(c) notes define
\begin{equation}
 K=\mathcal F^{-1}X\mathcal F,                           \tag{62}
\end{equation}
where $X$ is the position/scaling-coordinate operator.  Thus $K$ is a
first-order unbounded operator; it was not defined as a Kolmogorov factor of
the Hodge quadratic form.  Consequently the use of $Kf$ as though it were the
source energy feature in (56) requires the Gram identity (60), which is not a
formal consequence of differentiated Poisson.

In fact the identity fails on the full natural domain.

\begin{theorem}[Order mismatch no-go --- \PROVED]
Let $K=\mathcal F^{-1}X\mathcal F$ on $L^2(\mathbb R)$ and let
$G_\Gamma$ be the archimedean Weil operator whose Fourier multiplier has
logarithmic growth.  Then no isometry $J$ can satisfy
\begin{equation}
 J Kf=G_\Gamma^{1/2}f                              \tag{63}
\end{equation}
for all functions in a common invariant core.  In particular
$K^*K\ne G_\Gamma$ and the proposed observation (56) cannot hold with the
row-(c) operator $K$.
\end{theorem}

\begin{proof}
By unitarity of Fourier,
\[
 K^*K=\mathcal F^{-1}X^2\mathcal F.
\]
Choose normalized Fourier-localized functions $f_n$ supported where
$|x-n|\le1$.  Then
\[
 \langle K^*Kf_n,f_n\rangle=n^2+O(n).
\]
The Gamma multiplier obtained from the logarithmic derivative of the Gamma
factor is $O(\log(2+|x|))$; hence
\[
 \langle G_\Gamma f_n,f_n\rangle=O(\log n).
\]
The two quadratic forms cannot agree.  If (63) held with $J$ isometric, their
norm squares would agree for every $f_n$, a contradiction.
\end{proof}

The finite Euler terms do not repair the order mismatch: at a fixed threshold
they are a finite sum of bounded translations, hence contribute $O(1)$ on the
sequence above.  Therefore the full finite-threshold Hodge energy still cannot
be $K^*K$.

\paragraph{Consequence for G40.}
The unitary midpoint theorem (44)--(47) is correct, but its application
(48)--(51) with the row-(c) vector $Kf$ is \AUDIT{} and cannot be used as a
proof of G40.  Differentiated Poisson proves a connection/amplitude identity,
not equality of source and Hodge energy features.  A valid midpoint proof
would require an independently constructed feature $B_T$ satisfying
\[
 B_T^*B_T=S_M'
\]
and a source-defined unitary acting on its range whose difference port is the
G39 current.  Taking $B_T=(S_M')^{1/2}$ without deriving that unitary from the
Gamma--Tate network would merely restate the missing theorem.

\section{Unitary-orbit audit of the G39 auxiliary port}

For vectors/operators $W,Z$ in one Hilbert current space, a unitary $U$ can
satisfy
\begin{equation}
 Z=-\frac12(I-U)W                                      \tag{64}
\end{equation}
only if
\begin{equation}
 W^*Z+Z^*W+2Z^*Z=0.                                    \tag{65}
\end{equation}
Indeed $UW=W+2Z$, and (65) is exactly
$(W+2Z)^*(W+2Z)=W^*W$.

\begin{proposition}[Orthogonal-port no-go --- \PROVED]
If $Z\perp W$ and $Z\ne0$, no unitary $U$ satisfies (64).
\end{proposition}

\begin{proof}
Under orthogonality, (65) reduces to $2Z^*Z=0$, a contradiction.
\end{proof}

The G39 construction (27e) was adjoined as an auxiliary direct-sum port to the
Hodge current.  In that realization it is orthogonal to $W_{\Hod}$, so it
cannot be the midpoint residual asserted in (48).  A successful G40 network
must instead construct a nonorthogonal embedding
\[
 \jmath_T:\mathcal K_\infty^{(1)}\to\mathcal K_{\Hod,T}
\]
such that
\begin{align}
 H^*\jmath_TW_\infty&=-B_Q,\notag\\
 W_{\Hod}^*\jmath_TW_\infty
 +(\jmath_TW_\infty)^*W_{\Hod}
 +2(\jmath_TW_\infty)^*\jmath_TW_\infty&=0.            \tag{66}
\end{align}
The first line is amplitude; G39 proves it only before this embedding.  The
second line is the exact unitary-orbit energy constraint.  Constructing
$\jmath_T$ from Gamma--Tate, independently of (66), would give the desired
midpoint bridge.  Choosing it by solving (66) is equivalent to the missing
passivity theorem.

\section{Oriented-delay realization of the positive source metric}

Let $E_T:L^2(-T,T)\to L^2(\mathbb R)$ be extension by zero, let $S_u$ be
the unitary translation group on the \emph{ambient} space $L^2(\mathbb R)$,
and let $\mathfrak F$ be an involutive unitary satisfying
\begin{equation}
 \mathfrak F S_u\mathfrak F=S_{-u}.                     \tag{67}
\end{equation}
For $u>0$ define the oriented-delay feature
\begin{equation}
 \mathcal B_uF:=\frac1{\sqrt2}
 \binom{(I-S_u)F}{(I-S_{-u})F}.                         \tag{68}
\end{equation}

\begin{theorem}[Oriented-delay Fourier colligation --- \PROVED]
Define $\widehat{\mathfrak F}$ on the doubled current space by
\begin{equation}
 \widehat{\mathfrak F}\binom{x}{y}
 :=\binom{\mathfrak Fy}{\mathfrak Fx}.                 \tag{69}
\end{equation}
Then $\widehat{\mathfrak F}$ is an involutive unitary,
\begin{align}
 \mathcal B_u^*\mathcal B_uF
 &\text{ has quadratic form }\|(I-S_u)F\|^2,           \tag{70}\\
 \widehat{\mathfrak F}\mathcal B_u
 &=\mathcal B_u\mathfrak F.                            \tag{71}
\end{align}
Consequently the midpoint currents
\begin{equation}
 \mathcal B_{u,\pm}:=\frac12
 (I\pm\widehat{\mathfrak F})\mathcal B_u              \tag{72}
\end{equation}
satisfy the exact positive conservation law
\begin{equation}
 \mathcal B_u^*\mathcal B_u
 =\mathcal B_{u,+}^*\mathcal B_{u,+}
 +\mathcal B_{u,-}^*\mathcal B_{u,-}.                  \tag{73}
\end{equation}
\end{theorem}

\begin{proof}
Equation (69) is plainly an involutive unitary.  Since $S_u$ is unitary,
$\|(I-S_{-u})F\|=\|(I-S_u)F\|$, which proves (70).
From (67),
\[
 \mathfrak F(I-S_{-u})=(I-S_u)\mathfrak F,
 \qquad
 \mathfrak F(I-S_u)=(I-S_{-u})\mathfrak F,
\]
giving (71).  Equation (73) is the unitary midpoint identity.
\end{proof}

Let
\[
 w_\Gamma(u)=\frac{e^{-u/2}}{1-e^{-2u}}\qquad(u>0),
 \qquad c_n=\frac{\Lambda(n)}{\sqrt n}.
\]
The completed finite-threshold feature is the direct integral/sum
\begin{equation}
 \mathcal B_TF:=
 \left(
   \{w_\Gamma(u)^{1/2}\mathcal B_uF\}_{u>0},
   \{c_n^{1/2}\mathcal B_{\log n}F\}_{n<e^{2T}}
 \right).                                               \tag{74}
\end{equation}
Using the already proved Gamma delay identity, (70) gives
\begin{equation}
 \|\mathcal B_TF\|^2
 =\int_0^\infty w_\Gamma(u)\|(I-S_u)F\|^2du
 +\sum_{n<e^{2T}}c_n\|(I-S_{\log n})F\|^2.             \tag{75}
\end{equation}
Here the argument of every $\mathcal B_u$ is $E_TF$.  Thus $\mathcal B_T$
is a source-defined Kolmogorov factor of the positive Euler--Gamma reference
metric.  The direct-integral unitary obtained from (69)
intertwines $\mathcal B_T$ with $\mathfrak F$ and gives
\begin{equation}
 \mathcal B_T^*\mathcal B_T
 =\mathcal B_{T,+}^*\mathcal B_{T,+}
 +\mathcal B_{T,-}^*\mathcal B_{T,-}.                  \tag{76}
\end{equation}
This construction has the correct logarithmic spectral order and replaces the
invalid use of $K=\mathcal F^{-1}X\mathcal F$.

The preceding display proves a Gram identity for the positive reference.
It does \emph{not} identify that Gram with the independently defined Hodge
operator $\cL$ or with $S_M'$.  Defining $\cL:=\mathcal B_T^*\mathcal B_T$
would merely rename the reference form and would not prove compatibility
with the old factor used by the threshold induction.  The former equations
(76a)--(76b) are therefore withdrawn and classified as \AUDIT{}.

\paragraph{Interface still to prove.}
For G40 it remains to restrict the difference port to the harmonic newborn
input and verify the exact identity
\begin{equation}
 H^*\mathcal B_{T,-}Rv=B_Q.                              \tag{78}
\end{equation}
Equation (78) is the required identification of the
oriented difference port with the algebraic G39 port.  Unlike the former
$K$-bridge, (78) has the correct logarithmic order.  It remains only a
candidate until both its type map and its compatibility with the independently
defined Hodge/old Gram have been constructed.

\section{Authoritative balanced factorization of the supplied row (d)}

We now use the definitions in the supplied
\texttt{row-d-local-analysis.tex}.  Let $I_T=(-T,T)$, let $E_T$ denote
extension by zero to $\mathbb R$, and put
\[
 S_a:=E_T^*\tau_aE_T,
 \qquad (\tau_aF)(t)=F(t+a).
\]
Then $S_a^*=S_{-a}$.  For $a>0$ define the two \emph{ambient}
zero-extension channels from $L^2(I_T)$ to $L^2(\mathbb R)$
\begin{equation}
 J_{a,-}F:=\frac{\tau_aE_TF-E_TF}{\sqrt2},
 \qquad J_{a,+}F:=\frac{\tau_aE_TF+E_TF}{\sqrt2}.     \tag{79}
\end{equation}
They obey the exact polarization identity
\begin{equation}
 J_{a,-}^*J_{a,-}-J_{a,+}^*J_{a,+}
 =-(S_a+S_{-a}).                                      \tag{80}
\end{equation}
Indeed $\tau_a$ is unitary on the ambient line and
$E_T^*\tau_aE_T=S_a$.  This is the exact zero-extension factorization used
by the supplied row (d).  Both signs receive the same diagonal energy.

The Gamma multiplier in the supplied file is
\[
 g_\Gamma(\tau)=\Re\psi\left(\frac14+\frac{i\tau}{2}\right)
                  -\psi\left(\frac14\right).
\]
The integral formula for the digamma function and the substitution $t=2u$
give
\begin{equation}
 g_\Gamma(\tau)
 =2\int_0^\infty
 \frac{e^{-u/2}}{1-e^{-2u}}(1-\cos(\tau u))\,du.       \tag{81}
\end{equation}
Equivalently, on the natural form core,
\begin{equation}
 \langle G_{\Gamma,T}F,F\rangle
 =\int_0^\infty\frac{e^{-u/2}}{1-e^{-2u}}
       \|(I-\tau_u)E_TF\|_{L^2(\mathbb R)}^2\,du.     \tag{82}
\end{equation}
Thus $G_{\Gamma,T}\ge0$ and (82) is a source-defined Gamma Kolmogorov
factor, not a square root chosen from the desired signed operator.

Recall
\[
 m_0=\log\pi+\gamma+\frac\pi2+3\log2
     =\log\pi-\psi\left(\frac14\right)>0,
 \qquad w_n=\frac{\Lambda(n)}{\sqrt n}.
\]
Let $\mathcal X_T$ be the direct-sum analysis map whose components are the
Gamma feature in (82) and $w_n^{1/2}J_{\log n,-}$, and let
$\mathcal Y_T$ have components $m_0^{1/2}I$ and
$w_n^{1/2}J_{\log n,+}$, for $2\le n<e^{2T}$.

\begin{theorem}[Exact Gamma--Tate balanced factorization --- \PROVED]
On the form domain of the localized primitive operator in the supplied row
(d),
\begin{equation}
 A_T=\mathcal X_T^*\mathcal X_T-
     \mathcal Y_T^*\mathcal Y_T.                       \tag{83}
\end{equation}
After imposing the two Tate moments, with $P_T$ the orthogonal projection
onto $\mathcal P_T=\ker M_-\cap\ker M_+$,
\begin{equation}
 P_TA_TP_T=(\mathcal X_TP_T)^*(\mathcal X_TP_T)
           -(\mathcal Y_TP_T)^*(\mathcal Y_TP_T).      \tag{84}
\end{equation}
In particular this gives explicit, source-defined meanings to the abstract
factors $X_T,Y_T$ used in the return calculation of the supplied file.
\end{theorem}

\begin{proof}
The Gamma and constant pieces of (83) are respectively
$G_{\Gamma,T}$ and $-m_0I$.  For every prime power, (80), multiplied by
$w_n$, is exactly
$-w_n(S_{\log n}+S_{-\log n})$.  Summing gives the displayed definition of
$A_T$.  Compression by $P_T$ proves (84).
\end{proof}

The factorization also isolates the missing theorem without any notation
from v83--v86.

\begin{theorem}[Exact Douglas form of the remaining gate --- \PROVED]
The following are equivalent:
\begin{enumerate}
\item $P_TA_TP_T\ge0$;
\item $\|\mathcal Y_TP_TF\|\le\|\mathcal X_TP_TF\|$ for every $F$;
\item there is a contraction
\[
 C_T:\overline{\Ran(\mathcal X_TP_T)}
       \longrightarrow\overline{\Ran(\mathcal Y_TP_T)}
 \]
such that $C_T\mathcal X_TP_T=\mathcal Y_TP_T$.
\end{enumerate}
\end{theorem}

\begin{proof}
The equivalence of the first two statements is (84).  The inequality makes
$C_T(\mathcal X_TP_TF):=\mathcal Y_TP_TF$ well-defined and contractive on
the range; continuity extends it to the closure.  Conversely such a
contraction gives the norm inequality.
\end{proof}

There is a useful refinement: the source angular object exists before its
contractivity is known.

\begin{theorem}[Canonical source angular operator --- \PROVED]
The Gamma feature in $\mathcal X_T$ is injective.  Hence the rule
\begin{equation}
 \mathfrak C_T^{\rm src}(\mathcal X_TP_TF)
 :=\mathcal Y_TP_TF                                    \tag{84a}
\end{equation}
defines a linear operator on $\Ran(\mathcal X_TP_T)$ without assuming row
(d).  In the sense of sesquilinear forms on source vectors, its defect pulls
back to the primitive operator:
\begin{equation}
 \langle P_TA_TP_TF,G\rangle
 =\langle\mathcal X_TP_TF,\mathcal X_TP_TG\rangle
 -\langle\mathfrak C_T^{\rm src}\mathcal X_TP_TF,
          \mathfrak C_T^{\rm src}\mathcal X_TP_TG\rangle. \tag{84b}
\end{equation}
For every old/new decomposition of the primitive source, restriction of
(84a) gives the same old columns as (83); thus its compatibility with the
authoritative old factor is automatic.  It extends to a contraction on the
closed range if and only if row (d) holds at $T$.
\end{theorem}

\begin{proof}
By (81), $g_\Gamma(\tau)>0$ for every $\tau\ne0$.  If the Gamma feature of
$F$ vanishes, Plancherel gives
$g_\Gamma(\tau)|\widehat{E_TF}(\tau)|^2=0$ almost
everywhere.  An $L^2$ function supported at the singleton $\{0\}$ is zero,
so $E_TF=0$ and $F=0$.  Therefore $\mathcal X_T$ and its restriction to
$\mathcal P_T$ are injective, proving that (84a) is well defined.

For source vectors $F,G$, the sesquilinear form of the right side of (84b)
is
\[
 \langle\mathcal X_TP_TF,\mathcal X_TP_TG\rangle
 -\langle\mathcal Y_TP_TF,\mathcal Y_TP_TG\rangle,
\]
which is (84) by polarization.  Restriction to old/new columns commutes with
the source definition, proving compatibility.  Finally, bounded extension
with norm at most one is equivalent to the norm inequality in the preceding
theorem.
\end{proof}

Thus inverse engineering has produced a correctly typed, source-defined
bridge rather than a square root of the desired Schur remainder.  Its
unproved property is now a single quantitative statement,
\begin{equation}
 \boxed{\|\mathfrak C_T^{\rm src}\|\le1
        \quad\hbox{for every }T>0.}                    \tag{84c}
\end{equation}
Within the authoritative balanced factorization, no normalization or
old-column ambiguity remains in (84a)--(84b); the missing content is precisely
the global passivity bound (84c).  Identification of these columns with the
separately introduced v83--v86 Hodge factor remains subject to the Gram
compatibility audit (41).

\begin{theorem}[Correctly typed Poisson--Sonin observation --- \PROVED]
Let $U_T^X$ be any source-defined isometry from
$\overline{\Ran(\mathcal X_TP_T)}$ into a proposed normalized Sonin current
space.  Define on $U_T^X\Ran(\mathcal X_TP_T)$
\begin{equation}
 \Pi_T^{\rm src}\,U_T^X\mathcal X_TP_TF
 :=\mathcal Y_TP_TF.                                   \tag{86}
\end{equation}
Then $\Pi_T^{\rm src}$ is well defined, and
\begin{align}
 &\langle P_TA_TP_TF,G\rangle \notag\\
 &\quad=
 \langle U_T^X\mathcal X_TP_TF,U_T^X\mathcal X_TP_TG\rangle
 -\langle\Pi_T^{\rm src}U_T^X\mathcal X_TP_TF,
          \Pi_T^{\rm src}U_T^X\mathcal X_TP_TG\rangle. \tag{87}
\end{align}
The observation $\Pi_T^{\rm src}$ extends contractively to the closed current
space if and only if row (d) holds at $T$.  Requiring it to be isometric would
instead force $P_TA_TP_T=0$.
\end{theorem}

\begin{proof}
Injectivity of $\mathcal X_TP_T$ and isometry of $U_T^X$ make (86) well
defined.  Isometry and (84b) give (87).  The contractive-extension criterion
is exactly (84c).  If $\Pi_T^{\rm src}$ were isometric, the two terms on the
right of (87) would agree for every $F,G$, so the compressed primitive form
would vanish identically.
\end{proof}

Thus the normalized Sonin map $\mathcal I_T$ from (53) can serve as the
isometry $U_T^X$ only after one supplies a correctly typed source isometry
\begin{equation}
 V_T:\overline{\Ran(\mathcal X_TP_T)}\longrightarrow\mathcal S_T,
 \qquad V_T^*V_T=I,                                    \tag{88}
\end{equation}
derived from the Gamma--Tate Gram.  The actual normalized transport is then
$U_T^X=\mathcal I_TV_T$, and (86), not an isometric observation, is the
correct Poisson--Sonin--Gamma--Tate diagram.  Equation (88) is the precise
Gram/type interface absent from the supplied files.  Even after (88), the
contractivity of (86) remains the global sign theorem.

\paragraph{Audit against the Sonin metric in \texttt{main.tex}.}
The later section ``A derived metric: the Sonin compression'' defines instead
the positive Hilbert--Schmidt Gram
\[
 M(f)_{mn}=\operatorname{Tr}
 \bigl(\vartheta(\psi_m^f)\mathbf S
       \vartheta(\psi_n^f)^*\bigr).
\]
That positivity is unconditional, but the supplied manuscript explicitly
states that comparison with the Weil form is proved only in the prime-free
window and that the semilocal extension is not carried out.  For general
windows it leaves open the quantitative lower-frame estimate labelled
\texttt{eq:framebound}.  No equality in the supplied files identifies this
$M(f)$ with $(\mathcal X_TP_T)^*(\mathcal X_TP_T)$, so it does not construct
$V_T$ in (88).  Treating two positive Grams as identical merely because both
come from Sonin/Gamma data would be a Gram substitution and is forbidden.

\begin{theorem}[Global equivalence audit --- \PROVED]
The following assertions are equivalent:
\begin{enumerate}
\item $\|\mathfrak C_T^{\rm src}\|\le1$ for every $T>0$;
\item $P_TA_TP_T\ge0$ for every $T>0$;
\item the Weil quadratic form has the required sign on every compactly
supported test function satisfying the two Tate moment conditions;
\item the Riemann hypothesis holds.
\end{enumerate}
\end{theorem}

\begin{proof}
The equivalence of the first two assertions is (84b)--(84c).  Extension by
zero identifies the union over $T$ of the primitive spaces with the compactly
supported two-moment test core, so the second and third assertions are
equivalent.  The equivalence of the third and fourth assertions is Weil's
positivity criterion, with the sign convention fixed by
$A_T=-B_{\rm nuc}$ in the supplied manuscript.
\end{proof}

This theorem is not a proof of the final sign.  It records that any derivation
of (84c) from Poisson--Sonin data is itself a proof of RH and therefore must
contain a new global estimate beyond the currently available local
normalizations and intertwining identities.

\begin{theorem}[No channelwise Gamma--Tate contraction --- \PROVED]
Fix $a>0$.  For all sufficiently large $T$ there is an
$F\in\mathcal P_T$ such that
\begin{equation}
 \|J_{a,+}F\|^2-\|J_{a,-}F\|^2
 =2\Re\langle S_aF,F\rangle>0.                         \tag{89}
\end{equation}
Consequently a contraction implementing the global identity in the preceding
theorem cannot be a direct sum of contractions which send each arithmetic
difference channel $J_{a,-}$ to the corresponding sum channel $J_{a,+}$.
It must mix different places and/or the Gamma and Tate-terminal channels.
\end{theorem}

\begin{proof}
Choose a nonzero nonnegative smooth bump $\varphi$ supported in an interval
of length smaller than $a/4$.  For $T$ large, choose three numbers $x_j$,
separated from one another by more than $3a$, so that all six translates
supporting
\[
 F_j(t):=\varphi(t-x_j)+\varphi(t-x_j-a),\qquad 1\le j\le3,
\]
lie in $I_T$ and are mutually disjoint.  The two Tate moments of the three
$F_j$ form a $2\times3$ matrix.  Hence there is a nonzero real vector
$c=(c_1,c_2,c_3)$ such that
$F=\sum_jc_jF_j$ satisfies $M_-F=M_+F=0$.

By the disjointness of the three pairs, cross terms between different $j$
vanish.  Inside each pair exactly one overlap survives after translation by
$a$, and therefore
\[
 \Re\langle S_aF,F\rangle
 =\|\varphi\|^2\sum_{j=1}^3c_j^2>0.
\]
Expanding (79) gives (89).  A channelwise contraction would imply the
opposite weak inequality for every primitive $F$, a contradiction.
\end{proof}

This no-go result is important for the Poisson--Sonin proposal.  Polar
normalization of each local multiplier, even when isometric, cannot be the
missing $C_T$.  Any successful source construction must contain a genuinely
global off-diagonal mixing law whose Gram includes the centered atomic--
continuous cross $Q_c$.

\paragraph{Normalization audit.}
Equations (79)--(83) use all available local normalizations: the unique
self-dual coefficient $w_n=\Lambda(n)/\sqrt n$, the constant $m_0$, the
complete Gamma place, and the two exact Tate moments.  They determine the
two Grams and their signed difference, but they do not supply the contraction
$C_T$.  Producing $C_T$ from (83) by polar decomposition is legitimate only
after proving the norm inequality, and is therefore circular.  In the
threshold notation of the supplied file, the same missing assertion is the
factorization of the centered cross $Q_c$ through $D_0^{1/2}$ with unit
shorted-capacity budget.  Hence (83) is a genuine completion of the balanced
factorization bookkeeping, but not a proof of G40.

\paragraph{Notation audit.}
The symbol $\mathcal H_{5/4}$ occurs once in the supplied row (d), in the
definition of $\widehat{\mathcal R}_T$, and is not defined in any of the five
supplied files.  Formula (82) provides an explicit positive Gamma reference.
Identifying $\mathcal H_{5/4}$ with that reference requires either a stated
definition or a separately proved equality of forms; it must not be used
silently in the G40 argument.

The labels G39--G42 likewise do not occur in any of the five supplied files.
In this audit they refer only to the v83--v87 decomposition.  The authoritative
row-(d) target is (84), or equivalently the unit Douglas contraction above;
no status assigned to a G-label may override that operator inequality.

\section{Primitive-potential reduction and the continuous defect kernel}

The two Tate moments admit an exact differential parametrization which
removes the moving two-moment projection from the global estimate.  Put
\[
 L:=\partial_t^2-\frac14,
 \qquad H_0^2(I_T):=\{\phi\in H^2(I_T):
 \phi(\pm T)=\phi'(\pm T)=0\}.
\]

\begin{theorem}[Tate potential isomorphism --- \PROVED]
The map
\begin{equation}
 L:H_0^2(I_T)\longrightarrow
 \mathcal P_T=\ker M_-\cap\ker M_+                 \tag{90}
\end{equation}
is a continuous bijection onto the $L^2$ primitive space.  Its inverse is
the causal Green formula
\begin{equation}
 \phi(t)=\int_{-T}^{t}2\sinh\!\left(\frac{t-u}{2}\right)F(u)\,du.
                                                               \tag{91}
\end{equation}
After extension by zero, $L$ commutes distributionally with every ambient
translation $\tau_a$.
\end{theorem}

\begin{proof}
The kernel $k(x)=2\sinh(x/2)\mathbf1_{x\ge0}$ satisfies $Lk=\delta_0$.
Thus (91) gives $L\phi=F$ and the left endpoint conditions.  For $t\ge T$ it
equals
\[
 e^{t/2}\int_{-T}^{T}e^{-u/2}F(u)\,du
 -e^{-t/2}\int_{-T}^{T}e^{u/2}F(u)\,du,
\]
which vanishes together with its derivative exactly when $M_-F=M_+F=0$.
Hence its zero extension lies in $H^2(\mathbb R)$ and gives the right endpoint
conditions.  Conversely, integration by parts twice shows that $L\phi$ has
both moments zero for every clamped $\phi$.  The homogeneous equation
$L\phi=0$ and the two left boundary conditions give injectivity.  Standard
one-dimensional Green estimates give continuity of the inverse.  Constant-
coefficient differentiation commutes with translations on distributions.
\end{proof}

Let
\begin{equation}
 a_T(\tau):=g_\Gamma(\tau)-m_0
 -2\sum_{2\le n<e^{2T}}\frac{\Lambda(n)}{\sqrt n}
        \cos(\tau\log n).                              \tag{92}
\end{equation}

\begin{theorem}[Exact continuous defect representation --- \PROVED]
For $F=L\phi\in\mathcal P_T$,
\begin{align}
 \langle A_TF,F\rangle
 &=\frac1{2\pi}\int_{\mathbb R}
       a_T(\tau)(\tau^2+\tfrac14)^2
       |\widehat\phi(\tau)|^2\,d\tau                 \tag{93}\\
 &=\int_0^\infty w_\Gamma(u)
       \|(I-\tau_u)L\phi\|^2\,du
   -m_0\|L\phi\|^2 \notag\\
 &\quad-sum_{2\le n<e^{2T}}\frac{\Lambda(n)}{\sqrt n}
       \bigl(\langle\tau_{\log n}L\phi,L\phi\rangle
             +\langle\tau_{-\log n}L\phi,L\phi\rangle\bigr). \tag{94}
\end{align}
Here every function is extended by zero before translation.  Thus row (d)
is equivalent to nonnegativity of the explicit scalar kernel (93) on Fourier
transforms of clamped potentials; no moving projection or unspecified Hodge
operator remains.
\end{theorem}

\begin{proof}
On the ambient line, Plancherel turns $S_a+S_{-a}$ into multiplication by
$2\cos(a\tau)$, and (81)--(82) give the Gamma term.  Since
$\widehat{L\phi}(\tau)=-(\tau^2+1/4)\widehat\phi(\tau)$, this proves (93).
Expanding the same terms before Fourier transformation gives (94).
\end{proof}

Formula (93) is not a pointwise multiplier test: $a_T$ changes sign, while
$\widehat\phi$ belongs to the Paley--Wiener class forced by support and four
clamped boundary conditions.  Replacing that constrained class by arbitrary
$L^2$ Fourier data would destroy the problem.  Any new Poisson--Tate square
decomposition must therefore act on this Paley--Wiener subspace and include
cross-place terms; theorem (89) rules out a placewise decomposition.

\subsection{The global differentiated Poisson current}

On Meyer's common nuclear Fr\'echet core let $Z$ be the Zeta convolution
operator, let $\partial$ denote multiplication by $\log x$ (the derivation of
multiplicative convolution), and put
\begin{equation}
 C:=Z\partial(Z^{-1})=sum_{n\ge2}\Lambda(n)U_n.       \tag{95}
\end{equation}
Let $\mathscr Jf(x)=x^{-1}f(x^{-1})$ and let $\mathcal F$ be the self-dual
additive Fourier transform.

\begin{theorem}[Differentiated Poisson mixing identity --- \PROVED]
For every $f$ in the Poisson core
$\mathcal H_\cap=\{f:f(0)=\mathcal Ff(0)=0\}$,
\begin{align}
 Z\partial f+\mathscr JZ\partial\mathcal Ff
 &=CZf+\mathscr JCZ\mathcal Ff                         \tag{96}\\
 &=(C+\mathscr JC\mathscr J)Zf.                        \tag{97}
\end{align}
Thus all Euler places and both orientations occur in one source-defined
Poisson current before any positivity or spectral decomposition is invoked.
\end{theorem}

\begin{proof}
The derivation rule and $C=Z\partial(Z^{-1})$ give
$\partial Z=-CZ$.  Differentiate the Poisson identity
$Zf=\mathscr JZ\mathcal Ff$.  Since
$\partial\mathscr J=-\mathscr J\partial$, one obtains
\[
 -CZf+Z\partial f
 =\mathscr JCZ\mathcal Ff-\mathscr JZ\partial\mathcal Ff,
\]
which is (96).  Applying $\mathscr J$ to Poisson gives
$Z\mathcal Ff=\mathscr JZf$ and proves (97).
\end{proof}

Under the self-dual critical normalization, conjugation by $\mathscr J$
turns $U_n$ into the opposite weighted translation; consequently the right
side of (97) is the global source representative of the symmetric arithmetic
operator appearing in (92).  Unlike (27e), it acts on the full continuum
Poisson core and contains cross-place mixing through $Z$.

Define the unnormalized Poisson current
\begin{equation}
 \mathcal W_{\rm P}f:=Z\partial f+\mathscr JZ\partial\mathcal Ff. \tag{98}
\end{equation}
Equations (96)--(97) prove its divergence, but not its physical Gram.  The
next required identity is the continuum metric comparison
\begin{equation}
 \langle\mathcal W_{\rm P}f,
          (ZZ^*)^\dagger\mathcal W_{\rm P}f\rangle
 \stackrel{?}{=}
 \int_{\mathbb R}a_T(\tau)|\widehat f(\tau)|^2
       \frac{d\tau}{2\pi}+\text{explicit Tate terminals},       \tag{99}
\end{equation}
with all compressions and domains made precise.  Formula (99) is recorded as
\OPEN{}, not as an identity: $Z$ is a topological isomorphism on Meyer's
weighted space, not a bounded invertible operator on the physical critical
$L^2$ space, so $(ZZ^*)^\dagger$ cannot be imported without a closed-form
domain construction.  The restricted Sonin polar normalization (53) is the
available mechanism for such a construction, but equality with the
Gamma--Tate Gram remains to be derived.

\begin{proposition}[Poisson-current Gram homogeneity obstruction --- \PROVED]
The divergence identity (97) cannot be converted into the signed arithmetic
term in $A_T$ by a parameter-independent quadratic normalization of the
Euler current alone.  More precisely, scale the contact operator in the
right side of (97) and set
$\mathcal W_{\rm E}(\lambda):=\lambda(C+\mathscr JC\mathscr J)Zf$.
Its divergence is linear in $\lambda$, whereas every quadratic Gram
\begin{equation}
 \mathcal W_{\rm E}(\lambda)^*M\mathcal W_{\rm E}(\lambda) \tag{100}
\end{equation}
with $M$ independent of $\lambda$ is quadratic in $\lambda$.  Hence no
contact-scale-stable identity of this form can equal the affine family
$G_{\Gamma,T}-m_0I-\lambda(C+C^*)$.
\end{proposition}

\begin{proof}
By definition
$\mathcal W_{\rm E}(\lambda)=\lambda\mathcal W_{\rm E}(1)$.
Therefore (100) is multiplied by $\lambda^2$, while the arithmetic
piece of (92) is multiplied by $\lambda$.  Equality on an interval of
$\lambda$ is impossible unless the arithmetic operator vanishes.
\end{proof}

This homogeneity argument does not rule out an isolated numerical coincidence
at $\lambda=1$; it rules out deriving the desired linear form as a structural
Gram identity of the differentiated Euler current.  A valid construction
must use an affine background/cross term or, equivalently, a metric variation.

The correct linear object is the first variation of the Zeta metric.  For
the weighted Euler family
\[
 Z_t=\sum_{n\ge1}n^{-t/2}U_n,\qquad
 C_t=\sum_{n\ge2}\Lambda(n)n^{-t/2}U_n,\qquad
 G_t=Z_t^*Z_t,
\]
one has on the common core
\begin{align}
 Z_t'&=-\frac12Z_tC_t,                                 \tag{101}\\
 G_t'&=-\frac12(C_t^*G_t+G_tC_t),                     \tag{102}\\
 -2\frac d{dt}\|Z_tf\|^2
 &=\langle(C_t^*G_t+G_tC_t)f,f\rangle.                \tag{103}
\end{align}
These identities are algebraic consequences of the Euler logarithmic
derivative and require no sign assumption.

Thus the viable global theorem is not positivity of a current norm, but
monotonicity of a \emph{normalized} Poisson--Sonin metric after the complete
Gamma and Tate boundary corrections have been included.  Writing
$\widetilde G_t$ for that still-to-be-constructed physical normalized Gram,
the required source theorem has the form
\begin{equation}
 -2P_T\widetilde G_t'P_T
 =\mathcal R_t^*\mathcal R_t,
 \qquad t=1,                                           \tag{104}
\end{equation}
with $\widetilde G_t$ and $\mathcal R_t$ defined independently of the left
side.  Equation (104) would supply the missing global square.  The raw
$G_t=Z_t^*Z_t$ does not suffice: its derivative (102) has no automatic sign,
and restricted Sonin polar normalization makes its internal Gram constant,
moving the whole derivative into the cutoff/Tate embeddings.  Computing
those boundary derivatives is therefore the remaining load-bearing step.

\subsection{Exact motion and energy of the two Tate terminals}

The Tate part of that boundary motion can be computed explicitly.  On
$H_T=L^2(-T,T)$ define, for $t>0$,
\begin{equation}
 B_tF:=\binom{\displaystyle\int_{-T}^TF(x)e^{-tx/2}\,dx}
                 {\displaystyle\int_{-T}^TF(x)e^{tx/2}\,dx},
 \qquad P_t:=I-B_t^*K_t^{-1}B_t,                       \tag{105}
\end{equation}
where
\begin{equation}
 K_t:=B_tB_t^*
 =\begin{pmatrix}
 2\sinh(tT)/t&2T\\2T&2\sinh(tT)/t
 \end{pmatrix}.                                       \tag{106}
\end{equation}
For $t>0$ the eigenvalues of $K_t$ are
$2\sinh(tT)/t\pm2T>0$, so the inverse is ordinary, not generalized.

\begin{theorem}[Tate boundary velocity and terminal square --- \PROVED]
For $F\in\ker B_t$,
\begin{equation}
 P_t'F=-B_t^*K_t^{-1}B_t'F.                            \tag{107}
\end{equation}
This vector is the unique minimum-norm normal velocity $\nu_t(F)$ satisfying
the differentiated primitive constraint
\begin{equation}
 B_t\nu_t(F)+B_t'F=0.                                  \tag{108}
\end{equation}
Its energy is the explicit positive terminal form
\begin{equation}
 \boxed{
 \|\nu_t(F)\|^2
 =\langle B_t'F,K_t^{-1}B_t'F\rangle_{\mathbb C^2}.}   \tag{109}
\end{equation}
At $t=1$, $B_t'F$ consists exactly of the two logarithmic Tate moments
\begin{equation}
 \frac12\binom{-\int xF(x)e^{-x/2}\,dx}
                    {\int xF(x)e^{x/2}\,dx}.           \tag{110}
\end{equation}
\end{theorem}

\begin{proof}
Differentiate $P_t=I-B_t^*K_t^{-1}B_t$ and apply the result to a vector with
$B_tF=0$.  Every term containing $B_tF$ vanishes, leaving (107).  Applying
$B_t$ proves (108).  The normal space is $\Ran B_t^*$; hence any other
solution of (108) differs from (107) by an element of $\ker B_t$, orthogonal
to it.  Finally
\[
 \|B_t^*K_t^{-1}B_t'F\|^2
 =\langle B_t'F,K_t^{-1}B_tB_t^*K_t^{-1}B_t'F\rangle,
\]
which is (109).  Direct differentiation of (105) gives (110).
\end{proof}

Writing $a_t=2\sinh(tT)/t$, $b=2T$, and
$B_t'F=(u_-,u_+)^t$, the terminal square is completely scalar:
\begin{equation}
 \|\nu_t(F)\|^2
 =\frac{|u_-+u_+|^2}{2(a_t+b)}
  +\frac{|u_--u_+|^2}{2(a_t-b)}.                       \tag{111}
\end{equation}
This separates the even and odd Tate boundary modes and proves that no
undetermined normalization constant remains in the terminal contribution.

For the potential parametrization $F=L\phi$ of (90), integration by parts
and the clamped boundary conditions give
\begin{equation}
 \boxed{
 B_1'(L\phi)=\frac12
 \binom{\displaystyle\int_{-T}^T\phi(x)e^{-x/2}\,dx}
       {\displaystyle\int_{-T}^T\phi(x)e^{x/2}\,dx}.} \tag{111a}
\end{equation}
Indeed
$L(xe^{x/2})=e^{x/2}$ and
$L(xe^{-x/2})=-e^{-x/2}$.  Thus the two moving Tate terminals of the
primitive vector are exactly the two ordinary exponential moments of its
potential.  This gives a direct type-compatible link between the
Paley--Wiener reduction (93) and the terminal Gram (109).

\begin{theorem}[Kato transport of the primitive metric --- \PROVED]
Let $M_t=M_t^*$ be a differentiable ambient metric and set
\begin{equation}
 \Omega_t:=[P_t',P_t].                                 \tag{112}
\end{equation}
Then $\Omega_t^*=-\Omega_t$.  If $U_t$ solves
$U_t'=\Omega_tU_t$, $U_{t_0}=I$, it is unitary and transports the primitive
spaces:
\begin{equation}
 P_tU_t=U_tP_{t_0}.                                    \tag{113}
\end{equation}
On the fixed source $P_{t_0}H_T$, the transported compression
\begin{equation}
 \widehat M_t:=P_{t_0}U_t^*M_tU_tP_{t_0}              \tag{114}
\end{equation}
satisfies
\begin{equation}
 \widehat M_t'
 =P_{t_0}U_t^*P_t
   \bigl(M_t'+[M_t,\Omega_t]\bigr)
   P_tU_tP_{t_0}.                                      \tag{115}
\end{equation}
\end{theorem}

\begin{proof}
The commutator of two selfadjoint operators is skew-adjoint, so the evolution
is unitary.  Differentiating $P_t^2=P_t$ gives
$P_t'=P_tP_t'+P_t'P_t$ and $P_tP_t'P_t=0$.  These imply
$P_t'+[P_t,\Omega_t]=0$, which proves (113) by differentiation.  Finally,
differentiate (114), use $U_t'=\Omega_tU_t$ and
$(U_t^*)'=-U_t^*\Omega_t$, and insert (113); this gives (115).
\end{proof}

For the raw Zeta metric $M_t=G_t=Z_t^*Z_t$, equations (102) and (115) give
the completely explicit primitive derivative
\begin{equation}
 \widehat G_t'
 =P_{t_0}U_t^*P_t
 \left[-\frac12(C_t^*G_t+G_tC_t)+[G_t,\Omega_t]\right]
 P_tU_tP_{t_0}.                                       \tag{116}
\end{equation}
Here $P_t'$ and hence $\Omega_t$ are given by (105)--(107), so (116) has no
unknown connection.  The remaining issue is not differentiation but the
identification of the complete Gamma-corrected physical metric whose
derivative equals the kernel (93), followed by a proof that the resulting
operator in (115) is negative semidefinite.

\subsection{The completed Euler--Gamma logarithmic metric}

The ambient metric can be identified before taking its critical boundary
value.  Put
\begin{equation}
 \Lambda_{\mathbb Q}(s):=\pi^{-s/2}\Gamma(s/2)\zeta(s),
 \qquad s=\frac t2+i\tau,                              \tag{117}
\end{equation}
and, away from its zeros and poles,
$\mathfrak M_t(\tau):=|\Lambda_{\mathbb Q}(t/2+i\tau)|^2$.

\begin{theorem}[Completed logarithmic-metric identity --- \PROVED]
For $t>2$ one has the absolutely convergent identity
\begin{align}
 2\partial_t\log\mathfrak M_t(\tau)
 &=2\Re\frac{\Lambda_{\mathbb Q}'}
                  {\Lambda_{\mathbb Q}}(t/2+i\tau)    \tag{118}\\
 &=\Re\psi\left(\frac t4+\frac{i\tau}{2}\right)
   -\log\pi
   -2\sum_{n\ge2}\Lambda(n)n^{-t/2}
          \cos(\tau\log n).                           \tag{119}
\end{align}
At $t=1$, the Gamma part of (119) is exactly
$g_\Gamma(\tau)-m_0$.  When (119) is continued distributionally to $t=1$
and paired with $|\widehat F(\tau)|^2$ for
$F\in\mathcal P_T$, the polar terms at $s=0,1$ vanish by the two Tate
moments, and every arithmetic term with $\log n\ge2T$ pairs to zero by
support.  The resulting pairing is precisely
$\langle A_TF,F\rangle$ in (93).
\end{theorem}

\begin{proof}
Differentiate the logarithm of (117).  Since $ds/dt=1/2$, differentiating
$\log|\Lambda_{\mathbb Q}(s)|^2$ gives
$\Re(\Lambda_{\mathbb Q}'/\Lambda_{\mathbb Q})(s)$;
the extra factor two yields (118).  The logarithmic derivative of the Gamma
factor is
$-\tfrac12\log\pi+\tfrac12\psi(s/2)$, while
$\zeta'/\zeta(s)=-\sum_{n\ge2}\Lambda(n)n^{-s}$ for
$\Re s>1$.  Taking twice the real part proves (119).

At $t=1$,
$\Re\psi(1/4+i\tau/2)-\log\pi=g_\Gamma(\tau)-m_0$ by
the definitions in the supplied row (d).  The distributional continuation
is the standard explicit-formula continuation already realized in row (c).
The Mellin values at $0$ and $1$ are the two moments $M_-,M_+$, so their
polar contributions vanish.  Finally
$\langle\tau_{\log n}F,F\rangle=0$ when $\log n\ge2T$ (endpoint equality is
null), which truncates the arithmetic sum and gives (92)--(93).
\end{proof}

This theorem supplies the previously unspecified complete physical
Gamma--Euler metric variation.  It does not yet supply a positive Hilbert
boundary metric at $t=1$: $\mathfrak M_1$ vanishes at critical-line zeros,
and continuation from $t>2$ crosses possible zeros with real part greater
than $1/2$.  A proof of (104) must therefore construct a source boundary
space whose positive defect kernel includes exactly the residues encountered
in this continuation, rather than applying an inverse to
$\mathfrak M_1$ pointwise.

\begin{proposition}[Three-lines inverse-metric obstruction --- \PROVED]
Let $H$ be a Hilbert-valued holomorphic function in a vertical strip, with
the usual Hardy decay.  The function
\begin{equation}
 N_H(\sigma):=\int_{\mathbb R}|H(\sigma+i\tau)|^2\,d\tau \tag{120}
\end{equation}
is log-convex.  If it is symmetric about $\sigma=1/2$, its right derivative
at the center is nonnegative.  Applied to
$H(s)=\Lambda_{\mathbb Q}(s)\Phi(s)$, however, this controls a derivative
weighted by $|\Lambda_{\mathbb Q}|^2$ and does not imply (93).

To obtain the unweighted physical logarithmic derivative from (120), the
boundary datum would have to be divided by $\Lambda_{\mathbb Q}$ (up to an
outer factor of the same modulus).  Such a division is holomorphic across a
domain only if the datum cancels every zero there; requiring this for the
full primitive test space is equivalent to requiring the relevant domain to
be zero-free.  Therefore a bare three-lines/functional-equation argument
cannot prove row (d) without importing the missing zero information.
\end{proposition}

\begin{proof}
Log-convexity of Hardy $L^2$ norms is the Hilbert-valued three-lines theorem;
symmetry makes the center a minimum.  Differentiating
$|\Lambda_{\mathbb Q}\Phi|^2$ produces the logarithmic derivative multiplied
by $|\Lambda_{\mathbb Q}Phi|^2$, together with the derivative of $\Phi$.
The target (93) instead has the arbitrary Paley--Wiener weight
$|\widehat F|^2$.  Matching the weights requires boundary modulus
$|\Phi|=|\widehat F|/|\Lambda_{\mathbb Q}|$.  A holomorphic quotient with
this prescribed cancellation cannot exist for arbitrary $F$ at a zero of
$\Lambda_{\mathbb Q}$.  Hence the cancellation step assumes the missing
zero-free assertion.
\end{proof}

\subsection{Residue boundary space and its exact signature}

For $F\in\mathcal P_T$ use the entire Mellin--Fourier transform
\begin{equation}
 \Phi_F(z):=\int_{-T}^TF(t)e^{izt}\,dt,                \tag{121}
\end{equation}
and for a nontrivial zero $\rho=\beta+i\gamma$ put
$z_\rho=(\rho-1/2)/i=\gamma-i(\beta-1/2)$.  The involution
$\rho\mapsto1-\bar\rho$ becomes $z\mapsto\bar z$.

\begin{theorem}[Exact residue Krein boundary --- \PROVED]
After the two polar modes have been killed by $F\in\mathcal P_T$, the
distributional boundary value in (118)--(119) has the source realization
\begin{equation}
 \langle A_TF,F\rangle
 =\sum_{\mathcal O}\langle
     \mathcal E_{\mathcal O}F,
     J_{\mathcal O}\mathcal E_{\mathcal O}F\rangle,   \tag{122}
\end{equation}
where the sum is over functional-equation orbits of nontrivial zeros,
with multiplicity.  For a critical-line orbit $\rho=1/2+i\gamma$,
\begin{equation}
 \mathcal K_{\mathcal O}=\mathbb C,\qquad
 J_{\mathcal O}=1,\qquad
 \mathcal E_{\mathcal O}F=\Phi_F(\gamma).             \tag{123}
\end{equation}
For an off-line two-point orbit
$\{\rho,1-\bar\rho\}$,
\begin{equation}
 \mathcal K_{\mathcal O}=\mathbb C^2,\qquad
 J_{\mathcal O}=\begin{pmatrix}0&1\\1&0\end{pmatrix},\qquad
 \mathcal E_{\mathcal O}F=
 \binom{\Phi_F(z_\rho)}{\Phi_F(\bar z_\rho)}.         \tag{124}
\end{equation}
In the eigenbasis of $J_{\mathcal O}$ its contribution is
\begin{equation}
 \frac12|\Phi_F(z_\rho)+\Phi_F(\bar z_\rho)|^2
 -\frac12|\Phi_F(z_\rho)-\Phi_F(\bar z_\rho)|^2.     \tag{125}
\end{equation}
Thus the source residue boundary is canonically Krein; its metric is Hilbert
on every residue fiber if and only if every nontrivial zero lies on the
critical line.
\end{theorem}

\begin{proof}
The nuclear Lefschetz identity in the supplied row (c), equivalently the
contour shift of (118), identifies the primitive explicit form with the sum
of the autocorrelation transform at the nontrivial zeros.  The entire
continuation of the real-line identity
$\widehat{F*F^*}(\tau)=|\Phi_F(\tau)|^2$ is
\[
 \Phi_F(z)\overline{\Phi_F(\bar z)}.
\]
For a fixed critical zero this is $|\Phi_F(\gamma)|^2$, giving (123).  The
two members of an off-line functional-equation orbit contribute
\[
 \Phi_F(z)\overline{\Phi_F(\bar z)}
 +\Phi_F(\bar z)\overline{\Phi_F(z)},
\]
which is exactly the Gram in (124).  Diagonalizing the exchange matrix gives
(125).  Absolute convergence for compactly supported autocorrelations is the
convergence statement used by the supplied nuclear explicit formula.
\end{proof}

The negative vector in (125) is not an artifact of coordinates.  If
$z\ne\bar z$ and neither point is one of the two polar evaluation points,
the map
\begin{equation}
 F\longmapsto\bigl(M_-F,M_+F,Phi_F(z),\Phi_F(\bar z)\bigr) \tag{126}
\end{equation}
from $L^2(-T,T)$ onto $\mathbb C^4$ is surjective: its four representing
exponentials have distinct exponents and are linearly independent on every
nonempty interval.  Consequently the restricted evaluation
$F\in\mathcal P_T\mapsto(\Phi_F(z),\Phi_F(\bar z))$ is onto
$\mathbb C^2$.  Hence each off-line residue fiber contains a genuine negative
direction already on the primitive source.  A positive Hodge realization
must eliminate that fiber structurally; it cannot turn (124) positive by a
change of basis.

\subsection{A new global Hodge object: the unitarized Poisson quotient}

Let $\pi_0(u)$ denote the centrally normalized scaling action on the
source-defined quotient $\mathcal H_-^0=\mathcal H_-/Z\mathcal H_\cap$ of
the supplied row (c); central normalization means that the spectral
character attached to a zero $\rho$ is $e^{u(\rho-1/2)}$.

\begin{theorem}[Amenable unitarization Hodge bridge --- \PROVED]
Assume that $\mathcal H_-^0$ admits a Hilbert completion
$\mathscr H_{\rm P}$ satisfying the following source conditions:
\begin{enumerate}
\item the quotient map and the nuclear test-action of row (c) extend
continuously to a dense subspace of $\mathscr H_{\rm P}$;
\item $\pi_0(u)$ extends to a strongly continuous group on
$\mathscr H_{\rm P}$;
\item there is a source-derived constant $M<\infty$ such that
\begin{equation}
 \sup_{u\in\mathbb R}\|\pi_0(u)\|\le M.               \tag{127}
\end{equation}
\end{enumerate}
Then there is an equivalent positive Hilbert metric
$\langle\cdot,\cdot\rangle_{\rm Hod}$ on $\mathscr H_{\rm P}$ for which
every $\pi_0(u)$ is unitary.  Consequently the generator of the centered
scaling group is skew-adjoint, every transpose eigencharacter supplied by
row (c) has exponent in $i\mathbb R$, and all nontrivial zeros satisfy
$\Re\rho=1/2$.
\end{theorem}

\begin{proof}
Let $m$ be a translation-invariant mean on $L^\infty(\mathbb R)$ and define
\begin{equation}
 \langle x,y\rangle_{\rm Hod}
 :=m_u\,\langle\pi_0(u)x,\pi_0(u)y\rangle_{\mathscr H_{\rm P}}. \tag{128}
\end{equation}
Uniform boundedness gives the upper estimate
$\|x\|_{\rm Hod}\le M\|x\|$.  Since (127) also applies to $-u$,
$\|\pi_0(u)x\|\ge M^{-1}\|x\|$, giving the lower estimate
$\|x\|_{\rm Hod}\ge M^{-1}\|x\|$.  Thus (128) is a positive equivalent
Hilbert norm.  Translation invariance of $m$ makes the group unitary in this
metric.  Stone's theorem makes its generator skew-adjoint.  A character
$e^{u(\rho-1/2)}$ belonging to the transpose spectrum of a unitary group is
bounded for both signs of $u$, hence $\Re(\rho-1/2)=0$.
\end{proof}

The new object is (128): it is constructed without selecting spectral
subspaces and, once (127) is proved, supplies the missing positive Hodge
metric directly.  The load-bearing analytic theorem has therefore been
reduced to the source estimate (127), together with construction of the
Hilbert completion.  Sonin compression is a natural candidate norm, but the
supplied files only prove positivity of each compressed trace; they do not
prove that the centered scaling action descends to the quotient as a
uniformly bounded group.

There is an equivalent closed-range formulation.  If the image
$Z\mathcal H_\cap$ is a closed invariant subspace in a critical Hilbert
realization on which centered scaling is unitary, then the Hilbert quotient
is itself unitary and (127) holds with $M=1$.  Conversely, failure of closed
range is exactly where an off-line hyperbolic residue plane can enter the
nuclear quotient without contradicting unitarity of the ambient scaling
action.  Thus a source proof of critical closed range for the Poisson map is
a concrete alternative formulation of the missing Hodge theorem.

\begin{theorem}[Local Witt Hilbert completion and global Zeta obstruction
--- \PROVED]
Give the algebraic Witt orbit the Hilbert norm for which
$\{\phi_r\}_{r\ge1}$ is orthonormal, and give the Dirichlet coefficient
space the norm of $\ell^2(\mathbb N)$.  Then the row-(b) intertwiner
\begin{equation}
 J_2\phi_r=\delta_r                                      \tag{129}
\end{equation}
extends unitarily, and every $V_n$ (equivalently, convolution by $\delta_n$)
is an isometry.  However Meyer's global Zeta sum has no bounded extension to
this Hilbert completion, since
\begin{equation}
 Z\delta_1=\sum_{n\ge1}\delta_n\notin\ell^2(\mathbb N). \tag{130}
\end{equation}
Thus the positive Hilbert structure available in row (b) does not by itself
construct the Hilbert Poisson quotient required in (127).
\end{theorem}

\begin{proof}
Equation (129) maps one orthonormal basis to the other.  The rule
$V_n\phi_r=\phi_{nr}$ maps the basis injectively to an orthonormal subset,
so it is isometric; the same holds for
$\delta_r\mapsto\delta_{nr}$.  Formula (130) follows directly from the
definition $Z=\sum_nU_n$, and its coefficient sequence has infinite squared
norm.  A bounded operator must send $\delta_1$ to a Hilbert vector, proving
the obstruction.
\end{proof}

For $t>1$, the weighted cyclic vector
\begin{equation}
 Z_t\delta_1=\sum_{n\ge1}n^{-t/2}\delta_n             \tag{131}
\end{equation}
does lie in $\ell^2$, with norm squared $\zeta(t)$.  Its divergence as
$t\downarrow1$ is exactly the critical loss of the naive Hilbert completion.
Any successful completion must renormalize this pole while retaining the
two-dimensional Tate quotient and the nontrivial residue boundary; simply
discarding the divergent cyclic direction would also discard the Poisson
extension data.

\begin{theorem}[Critical gcd boundary state --- \PROVED]
For $t>1$ put
\begin{equation}
 e_t:=\zeta(t)^{-1/2}\sum_{n\ge1}n^{-t/2}\delta_n
 \in\ell^2(\mathbb N).                                 \tag{132}
\end{equation}
Then, for all $m,n\ge1$ and $g=(m,n)$,
\begin{equation}
 \langle V_ne_t,V_me_t\rangle
 =\left(\frac{mn}{g^2}\right)^{-t/2}.                 \tag{133}
\end{equation}
Consequently the critical kernel
\begin{equation}
 K(m,n):=\lim_{t\downarrow1}\langle V_ne_t,V_me_t\rangle
 =\frac{(m,n)}{\sqrt{mn}}                              \tag{134}
\end{equation}
is positive definite.  It extends to the group positive-definite function
on $\mathbb Q_+^\times$
\begin{equation}
 \varphi(q)=\prod_p p^{-|v_p(q)|/2}.                   \tag{135}
\end{equation}
Its GNS representation is a unitary representation of
$\mathbb Q_+^\times$, carries a cyclic vector $\Omega$, and satisfies
$\langle U_n\Omega,U_m\Omega\rangle=K(m,n)$.
\end{theorem}

\begin{proof}
The supports of $V_ne_t$ and $V_me_t$ meet when $nr=ms$.  Writing
$m=ga$, $n=gb$ with $(a,b)=1$, all solutions are $r=ak$, $s=bk$.  Therefore
\[
 \langle V_ne_t,V_me_t\rangle
 =\frac1{\zeta(t)}\sum_{k\ge1}(ak)^{-t/2}(bk)^{-t/2}
 =(ab)^{-t/2},
\]
which is (133) and tends to (134).  A pointwise limit of positive-definite
kernels is positive definite.  Since
$K(m,n)=\varphi(m^{-1}n)$ and (135) is the corresponding reduced-prime
factorization, the GNS theorem gives the unitary group representation.
\end{proof}

This boundary is fully explicit.  For each prime let $r_p=p^{-1/2}$ and put
\begin{equation}
 d\mu_p(e^{i\theta})
 :=\frac{1-r_p^2}{|1-r_pe^{i\theta}|^2}\frac{d\theta}{2\pi}. \tag{136}
\end{equation}
Its Fourier moments are
$\int e^{ik\theta}d\mu_p=r_p^{|k|}$.  Hence the GNS space in the theorem is
realized by the infinite product probability space
\begin{equation}
 \mathscr H_{\gcd}:=L^2\left(\prod_p\mathbb T,
                              \bigotimes_p\mu_p\right), \tag{137}
\end{equation}
with $U_q$ multiplication by
$\prod_pz_p^{v_p(q)}$ and $\Omega=1$.  Inversion
$q\mapsto q^{-1}$ is unitary because every $\mu_p$ is invariant under
$z_p\mapsto\bar z_p$.

The gcd boundary supplies exactly the cross-place structure absent from a
direct sum of local ports and realizes the central coefficient
$\langle U_n\Omega,\Omega\rangle=n^{-1/2}$.  It does not yet include the
archimedean Gamma channel or prove that the physical observation from
clamped potentials into $\mathscr H_{\gcd}$ has Gram (93).  That observation
is now a concrete map to construct, rather than an unspecified positive
space.

\begin{theorem}[Gcd--Gamma reference observation --- \PROVED]
Let $\mathcal A_T=\{n:2\le n<e^{2T},\ \Lambda(n)>0\}$ and let
$K_T=(K(m,n))_{m,n\in\mathcal A_T}$.  This matrix is strictly positive.
Writing $\chi_n=U_n\Omega\in\mathscr H_{\gcd}$, define
\begin{equation}
 \eta_n:=\sum_{m\in\mathcal A_T}
          (K_T^{-1/2})_{mn}\chi_m.                    \tag{138}
\end{equation}
Then $\{\eta_n:n\in\mathcal A_T\}$ is orthonormal.  If
$\mathscr K_{\Gamma,T}$ denotes the Gamma delay space in (82), the map
\begin{equation}
 \mathcal V_T^X(\mathcal X_TF):=
 \left(
   \{w_\Gamma(u)^{1/2}(I-\tau_u)E_TF\}_{u>0},
   \sum_{n\in\mathcal A_T}\eta_n\otimes
        w_n^{1/2}J_{\log n,-}F
 \right)                                               \tag{139}
\end{equation}
is a source-defined isometry from
$\overline{\Ran\mathcal X_T}$ onto its generated range in
\begin{equation}
 \mathscr K_T^X:=\mathscr K_{\Gamma,T}
 \oplus(\mathscr H_{\gcd}\otimes L^2(\mathbb R)).      \tag{140}
\end{equation}
In particular
\begin{equation}
 (\mathcal V_T^X\mathcal X_T)^*
 (\mathcal V_T^X\mathcal X_T)=\mathcal X_T^*\mathcal X_T. \tag{141}
\end{equation}
\end{theorem}

\begin{proof}
The vectors $\chi_n$ are distinct characters on the product torus (137).
Because the product Poisson measure has full support, a finite linear
combination of them which vanishes in $L^2$ is the zero trigonometric
polynomial.  Thus their Gram $K_T$ is strictly positive.  Equation (138)
gives
\[
 (\langle\eta_m,\eta_n\rangle)_{m,n}
 =K_T^{-1/2}K_TK_T^{-1/2}=I.
\]
Orthogonality makes the squared norm of the arithmetic component of (139)
equal to
$\sum_nw_n\|J_{\log n,-}F\|^2$.  Equation (82) gives equality for the Gamma
component.  Their direct sum is precisely the defining Gram of
$\mathcal X_T$, proving (141).
\end{proof}

This closes the Gram/type interface (88) after replacing the previously
unspecified Sonin target by the new source-derived gcd--Gamma Hodge reference
space (140).  It is compatible with restriction to old contact sets because
all finite Grams are compressions of the single global kernel (134), although
the positive square roots $K_T^{-1/2}$ themselves need not be nested.  The
canonical object is the unorthogonalized GNS span; (138) is only a coordinate
choice inside it.

The opposite observation is defined on the same arithmetic labels by
replacing $J_{\log n,-}$ with $J_{\log n,+}$ and adjoining the constant
channel $m_0^{1/2}F$.  The resulting angular map is exactly
$\mathfrak C_T^{\rm src}$ transported by $\mathcal V_T^X$.  Its
contractivity is still (84c); the GNS construction proves that its domain
Gram and cross-place type are correct, but does not yet prove the norm-one
bound.

\begin{theorem}[Global gcd-oriented midpoint colligation --- \PROVED]
Let $I_T=\mathcal A_T\times\{+,-\}$ and put
$\chi_{n,+}=U_n\Omega$, $\chi_{n,-}=U_{n^{-1}}\Omega$.  Let
$\mathbf K_T$ be their Gram matrix and define the orthonormal family
\begin{equation}
 \eta_{n,\epsilon}:=\sum_{(m,\delta)\in I_T}
 (\mathbf K_T^{-1/2})_{(m,\delta),(n,\epsilon)}
 \chi_{m,\delta}.                                      \tag{142}
\end{equation}
The inversion unitary $\mathfrak J_{\gcd}$ satisfies
\begin{equation}
 \mathfrak J_{\gcd}\eta_{n,+}=\eta_{n,-},
 \qquad \mathfrak J_{\gcd}\eta_{n,-}=\eta_{n,+}.     \tag{143}
\end{equation}
Let $\mathfrak F$ be any involutive unitary on $L^2(\mathbb R)$ with
$\mathfrak F\tau_a\mathfrak F=\tau_{-a}$, and define
\begin{equation}
 \mathbb U_T:=\mathfrak J_{\gcd}\otimes\mathfrak F.   \tag{144}
\end{equation}
The arithmetic feature
\begin{equation}
 \mathbb B_T^{\rm ar}F
 :=\sum_{n\in\mathcal A_T}\frac{w_n^{1/2}}2
 \left[
  \eta_{n,+}\otimes(I-\tau_{\log n})E_TF
 +\eta_{n,-}\otimes(I-\tau_{-\log n})E_TF
 \right]                                               \tag{145}
\end{equation}
satisfies
\begin{align}
 (\mathbb B_T^{\rm ar})^*\mathbb B_T^{\rm ar}
 &=\sum_{n\in\mathcal A_T}w_nJ_{\log n,-}^*J_{\log n,-}, \tag{146}\\
 \mathbb U_T\mathbb B_T^{\rm ar}
 &=\mathbb B_T^{\rm ar}\mathfrak F.                  \tag{147}
\end{align}
Consequently
\begin{equation}
 \mathbb B_T^{\rm ar}
 =\mathbb B_{T,+}^{\rm ar}+\mathbb B_{T,-}^{\rm ar},
 \qquad
 (\mathbb B_T^{\rm ar})^*\mathbb B_T^{\rm ar}
 =\sum_{\epsilon=\pm}
   (\mathbb B_{T,\epsilon}^{\rm ar})^*
    \mathbb B_{T,\epsilon}^{\rm ar},                  \tag{148}
\end{equation}
where
$\mathbb B_{T,\pm}^{\rm ar}=\tfrac12(I\pm\mathbb U_T)
\mathbb B_T^{\rm ar}$.
\end{theorem}

\begin{proof}
The characters indexed by $I_T$ are distinct, so the proof used for (138)
shows $\mathbf K_T>0$.  Inversion permutes the characters by the swap
$P(n,+)=(n,-)$, hence $P\mathbf K_TP=\mathbf K_T$.
Functional calculus gives $P\mathbf K_T^{-1/2}P=\mathbf K_T^{-1/2}$,
proving (143).  Orthogonality of the $\eta_{n,\epsilon}$ and equality of the
two translation-difference norms give
\[
 \|\mathbb B_T^{\rm ar}F\|^2
 =\sum_n\frac{w_n}{4}
  \left(\|(I-\tau_{\log n})E_TF\|^2
       +\|(I-\tau_{-\log n})E_TF\|^2\right)
 =\sum_nw_n\|J_{\log n,-}F\|^2,
\]
which is (146).  Equations (143)--(144) and
$\mathfrak F(I-\tau_a)=(I-\tau_{-a})\mathfrak F$ give (147).
The unitary midpoint identity proves (148).
\end{proof}

Adding the Gamma-oriented feature (68), integrated with weight
$w_\Gamma(u)$, gives a single source-defined reference current
$\mathbb B_T^X$ with
\begin{equation}
 (\mathbb B_T^X)^*\mathbb B_T^X=\mathcal X_T^*\mathcal X_T, \tag{149}
\end{equation}
and an exact positive midpoint conservation law.  Therefore normalization,
global cross-place typing, reference Gram and conservation are closed in this
new gcd--Gamma colligation.  The remaining load-bearing identity is to show
that its appropriate midpoint observation equals the opposite factor
$\mathcal Y_T$ (including the constant and Tate-terminal channels) with no
extra energy.  That equality would prove (84c).

\begin{theorem}[M\"obius whitening of the gcd boundary --- \PROVED]
For a finite set of primes $S$, let $m_S$ be Haar probability measure on
$\mathbb T^S$, let $r_p=p^{-1/2}$, and put
\begin{equation}
 d\mu_S(z)=\prod_{p\in S}
 \frac{1-r_p^2}{|1-r_pz_p|^2}\,dm_S(z),
 \qquad
 b_S(z)=\prod_{p\in S}
 \frac{1-r_pz_p}{\sqrt{1-r_p^2}}.                     \tag{150}
\end{equation}
Then multiplication by $b_S$ is a unitary
\begin{equation}
 M_{b_S}:L^2(\mathbb T^S,m_S)longrightarrow
          L^2(\mathbb T^S,\mu_S),                     \tag{151}
\end{equation}
and
\begin{equation}
 |b_S(z)|^2d\mu_S(z)=dm_S(z).                         \tag{152}
\end{equation}
The numerator of $b_S$ is the finite M\"obius Euler multiplier and its
denominator is the unique positive normalization making it isometric.
\end{theorem}

\begin{proof}
For each prime the Poisson density in (136) gives
\[
 \frac{|1-r_pz_p|^2}{1-r_p^2}
 \frac{1-r_p^2}{|1-r_pz_p|^2}=1.
\]
Multiplying these identities proves (152), and integration proves (151).
\end{proof}

For the deformation $r_p(t)=p^{-t/2}$ the same construction gives
$b_{S,t}$.  Its unnormalized numerator is exactly $a_{S,t}$ in (20), while
the derivative of the positive denominator is real scalar.  Passing to the
polar phase therefore yields
\begin{equation}
 (M_{b_{S,t}}^{\rm polar})^*
 \partial_tM_{b_{S,t}}^{\rm polar}
 =\frac14M_{c_{S,t}-\overline{c_{S,t}}},               \tag{153}
\end{equation}
which agrees with (26).  Thus the gcd boundary, finite Euler polar
normalization and CCM Sonin normalization are three realizations of the same
source isometry at finite $S$.

What survives after whitening is not an Euler norm defect---that Hermitian
part is exactly absorbed by the change from $m_S$ to $\mu_S$---but the motion
of the Poisson/cutoff observation inside the whitened space.  This confirms
from a second construction that the remaining G40 term must be a boundary
second fundamental form, now combining (107)--(116) with the finite-product
measure variation in (150).

\begin{theorem}[Centered Euler score and exact Fisher budget --- \PROVED]
Let $r_p(t)=p^{-t/2}$ and let $\mu_{p,t}$ be the Poisson measure (136) with
$r_p$ replaced by $r_p(t)$.  Its logarithmic score is
\begin{equation}
 \sigma_{p,t}(z):=\partial_t\log\frac{d\mu_{p,t}}{dm}(z)
 =(\log p)\left[
 \frac{r_p(t)^2}{1-r_p(t)^2}
 -\Re\frac{r_p(t)z}{1-r_p(t)z}
 \right].                                              \tag{154}
\end{equation}
It is centered and has the exact Fisher energy
\begin{equation}
 \int_{\mathbb T}\sigma_{p,t}\,d\mu_{p,t}=0,
 \qquad
 \int_{\mathbb T}|\sigma_{p,t}|^2d\mu_{p,t}
 =\frac{(\log p)^2r_p(t)^2}
        {2(1-r_p(t)^2)^2}.                             \tag{155}
\end{equation}
For a finite set $S$, the product score
$\sigma_{S,t}=\sum_{p\in S}\sigma_{p,t}$ satisfies
\begin{equation}
 \|\sigma_{S,t}\|_{L^2(\mu_{S,t})}^2
 =\sum_{p\in S}\frac{(\log p)^2r_p(t)^2}
        {2(1-r_p(t)^2)^2}.                             \tag{156}
\end{equation}
Moreover, for every differentiable observable $H_t$,
\begin{equation}
 \partial_t\int H_t\,d\mu_{S,t}
 =\int\partial_tH_t\,d\mu_{S,t}
  +\int(H_t-\mathbb E_tH_t)\sigma_{S,t}\,d\mu_{S,t}. \tag{157}
\end{equation}
\end{theorem}

\begin{proof}
Differentiate the density in (150), using
$r_p'=-(\log p)r_p/2$, to obtain (154).  The Poisson moments are
$\int z^k d\mu_{p,t}=r_p(t)^k$ for $k\ge0$, so the expression is centered.
Put $r=r_p(t)$ and $W=rz/(1-rz)=\sum_{k\ge1}r^kz^k$.
The same moments give
\[
 \mathbb EW=\frac{r^2}{1-r^2},\qquad
 \operatorname{Var}(\Re W)=\frac{r^2}{2(1-r^2)^2}.
\]
This proves (155).  Independence makes distinct prime scores orthogonal and
gives (156).  Differentiation under the finite product integral gives (157);
centering the second term uses $\mathbb E_t\sigma_{S,t}=0$.
\end{proof}

Expanding (154),
\begin{equation}
 \sigma_{p,t}(z)
 =(\log p)\sum_{k\ge1}r_p(t)^k
       \bigl(r_p(t)^k-\Re z^k\bigr),                  \tag{158}
\end{equation}
so it contains every prime power with the central coefficient and subtracts
exactly its Poisson expectation.  Formula (157) is therefore a probability-
space realization of the centered atomic cross denoted $Q_c$ in the supplied
threshold analysis.  Cauchy--Schwarz gives the sharp covariance estimate
\begin{equation}
 |\operatorname{Cov}_t(H,\sigma_{S,t})|^2
 \le\operatorname{Var}_t(H)\,
      \|\sigma_{S,t}\|_2^2.                           \tag{159}
\end{equation}
To finish G40 by this route one must identify the old defect form $D_0$ with
the relevant conditional variance and show that the born/reference
normalization makes the Fisher factor in (159) exactly the available unit
budget.  The score and its complete prime-power budget are now explicit;
that final variance identification remains \OPEN{}.

\begin{theorem}[Exact arithmetic innovation budget --- \PROVED]
Order the finite contact characters by their integer label.  Let
$\mathscr E_{N-1}$ be the span of the characters $\chi_m=U_m\Omega$ with
$\Lambda(m)>0$ and $m<N$, and suppose $\Lambda(N)>0$.  Write $K_{N-1}$ for
their Gram matrix and
\begin{equation}
 k_N=(\langle\chi_m,\chi_N\rangle)_{m<N},
 \qquad
 \kappa_N:=1-k_N^*K_{N-1}^{-1}k_N.                    \tag{160}
\end{equation}
Then
\begin{equation}
 0<\kappa_N\le1,
 \qquad
 \|P_{\mathscr E_{N-1}}\chi_N\|^2
 =k_N^*K_{N-1}^{-1}k_N\le1,                           \tag{161}
\end{equation}
and the normalized innovation
\begin{equation}
 r_N:=\kappa_N^{-1/2}
 \left(\chi_N-sum_{m<N}(K_{N-1}^{-1}k_N)_m\chi_m\right) \tag{162}
\end{equation}
is a unit vector orthogonal to $\mathscr E_{N-1}$.
\end{theorem}

\begin{proof}
The enlarged character family is linearly independent in the full-support
product measure, so its Gram matrix
\[
 \begin{pmatrix}K_{N-1}&k_N\\k_N^*&1\end{pmatrix}
\]
is strictly positive.  Its Schur complement is $\kappa_N>0$; since the
subtracted term is nonnegative, $\kappa_N\le1$.  The normal equations for
orthogonal projection give the coefficients $K_{N-1}^{-1}k_N$, and direct
expansion gives (161)--(162).
\end{proof}

This theorem pays the exact unit Schur budget for every purely arithmetic
birth in the global gcd metric, with all prime-power correlations retained.
It strengthens mere all-depth summability: no estimate is used.  To transfer
(160)--(162) to the physical threshold capacity (41i), one must prove that
the Gamma--continuous and Tate-terminal columns enlarge
$\mathscr E_{N-1}$ so that the physical centered load $b_E$ is precisely the
projection load in this augmented Hilbert space.  That is the remaining
mixed observation identity.

\begin{corollary}[Closed prime-power innovation --- \PROVED]
Augment the old contact space by the Tate cyclic vector,
\begin{equation}
 \mathscr E_{N-1}^{\rm aug}
 :=\operatorname{span}\bigl(\Omega,
       \{\chi_m:\Lambda(m)>0,\ m<N\}\bigr).           \tag{163}
\end{equation}
If $N=p^k$, with the convention $\chi_{p^0}=\Omega$, then
\begin{align}
 P_{\mathscr E_{N-1}^{\rm aug}}\chi_{p^k}
 &=p^{-1/2}\chi_{p^{k-1}},                             \tag{164}\\
 \left\|\chi_{p^k}-
 P_{\mathscr E_{N-1}^{\rm aug}}\chi_{p^k}\right\|^2
 &=1-p^{-1}.                                           \tag{165}
\end{align}
The normalized innovations
\begin{equation}
 e_{p,k}(z):=\frac{(z_p-p^{-1/2})z_p^{k-1}}
                   {\sqrt{1-p^{-1}}}                  \tag{166}
\end{equation}
are orthonormal as $(p,k)$ varies.
\end{corollary}

\begin{proof}
For one Poisson factor with parameter $r=p^{-1/2}$,
\begin{equation}
 |z-r|^2d\mu_p(z)=(1-r^2)dm(z).                        \tag{167}
\end{equation}
Therefore the functions
$(z-r)z^{k-1}/\sqrt{1-r^2}$ are orthonormal and have mean zero.  Independence
makes functions belonging to different primes orthogonal.  Since
\[
 z_p^k-rz_p^{k-1}=(z_p-r)z_p^{k-1},
\]
the right side is orthogonal to the augmented old span, proving
(164)--(166).
\end{proof}

Thus the entire prime-power return chain is an exact AR(1)-type orthogonal
innovation process in the gcd boundary.  The factor $1-p^{-1}$ is the same
local denominator appearing in the M\"obius whitening (150).  This provides
an inverse-free, all-depth arithmetic capacity factorization; no Witt-moment
majorant is needed for the arithmetic part.

\begin{theorem}[Archimedean Gamma GNS boundary --- \PROVED]
Put $a_j=2j+\tfrac12$ and let
\begin{equation}
 d\nu_{a_j}(x):=\frac{a_j}{2}e^{-a_j|x|}\,dx
 \qquad(x\in\mathbb R).                                \tag{168}
\end{equation}
On $L^2(\nu_{a_j})$ let $U_\tau$ be multiplication by $e^{i\tau x}$ and
$\Omega_j=1$.  Then
\begin{align}
 \langle U_\tau\Omega_j,\Omega_j\rangle
 &=\frac{a_j^2}{a_j^2+\tau^2},                         \tag{169}\\
 \frac1{a_j}\|(U_\tau-I)\Omega_j\|^2
 &=\frac{2\tau^2}{a_j(a_j^2+\tau^2)}.                 \tag{170}
\end{align}
Consequently
\begin{equation}
 g_\Gamma(\tau)
 =\sum_{j\ge0}\frac1{a_j}\|(U_\tau-I)\Omega_j\|^2,  \tag{171}
\end{equation}
and the infinite product boundary has characteristic function
\begin{equation}
 \prod_{j\ge0}\frac{a_j^2}{a_j^2+\tau^2}
 =\frac{|\Gamma(\frac14+\frac{i\tau}{2})|^2}
        {\Gamma(\frac14)^2}.                           \tag{172}
\end{equation}
\end{theorem}

\begin{proof}
The Fourier transform of the two-sided exponential density (168) is
$a_j^2/(a_j^2+\tau^2)$, proving (169); expanding the difference norm proves
(170).  The digamma series with
$a_j=2(j+1/4)$ is exactly (171).  Finally the Weierstrass product for the
ratio
$\Gamma(1/4+i\tau/2)\Gamma(1/4-i\tau/2)/\Gamma(1/4)^2$
is the product in (172).
\end{proof}

For the metric deformation set $a_j(t)=2j+t/2$ and let $\nu_{a_j(t)}$ be
as in (168).  Its centered score and Fisher energy are
\begin{equation}
 \sigma_{\Gamma,j,t}(x)
 =\frac12\left(\frac1{a_j(t)}-|x|\right),
 \qquad
 \|\sigma_{\Gamma,j,t}\|_2^2=\frac1{4a_j(t)^2}.       \tag{173}
\end{equation}
Indeed $|x|$ is exponentially distributed with rate $a_j(t)$.  The scores
are independent in the product boundary and
$\sum_j1/(4a_j(t)^2)<\infty$, so the complete Gamma score is an honest
$L^2$ random variable.  Combining it with the finite Euler score (154)
gives a single centered Gamma--Euler score whose Fisher energies add.

The constant $m_0$ is not an external correction to this boundary.  It is
the scale coordinate paired with the Gamma ratio coordinate, as follows.

\begin{theorem}[Gamma ratio--scale normalization --- \PROVED]
Let $R_+,R_-$ be independent Gamma random variables with shape
$a=\tfrac14$ and unit scale, and put
\begin{equation}
 X:=\frac12\log\frac{R_+}{R_-},
 \qquad
 Y:=\log\pi-\frac12\log(R_+R_-).                    \tag{174}
\end{equation}
Then
\begin{align}
 \phi_\Gamma(\tau):=\mathbb E e^{i\tau X}
 &=\frac{|\Gamma(\frac14+\frac{i\tau}{2})|^2}
         {\Gamma(\frac14)^2},                       \tag{175}\\
 \mathbb EY&=\log\pi-\psi\left(\frac14\right)=m_0,\tag{176}\\
 \mathbb E\!\left(Ye^{i\tau X}\right)
 &=\bigl(m_0-g_\Gamma(\tau)\bigr)\phi_\Gamma(\tau).
                                                               \tag{177}
\end{align}
Moreover the conditional scale is the explicit positive function
\begin{equation}
 h_\Gamma(x):=\mathbb E(Y\mid X=x)
 =\log\pi-\psi\left(\frac12\right)+\log(2\cosh x)>0. \tag{178}
\end{equation}
Thus the finite positive measure $h_\Gamma\,d\mu_X$ has Fourier transform
$\phi_\Gamma(\tau)(m_0-g_\Gamma(\tau))$.
The ratio law itself is explicit:
\begin{equation}
 d\mu_X(x)=\frac{\sqrt2}{B(\frac14,\frac14)}
             (\cosh x)^{-1/2}\,dx.                  \tag{178a}
\end{equation}
\end{theorem}

\begin{proof}
For a unit-scale Gamma variable $R$ of shape $a$,
$\mathbb E R^z=\Gamma(a+z)/\Gamma(a)$.  Independence with
$z=\pm i\tau/2$ proves (175), and
$\mathbb E\log R=\psi(a)$ proves (176).  Differentiating the two Mellin
moments gives
\[
 \mathbb E\!\left(e^{i\tau X}\log R_+\right)
 =\phi_\Gamma(\tau)\psi\left(a+\frac{i\tau}{2}\right),
 \quad
 \mathbb E\!\left(e^{i\tau X}\log R_-\right)
 =\phi_\Gamma(\tau)\psi\left(a-\frac{i\tau}{2}\right).
\]
Substitution in (174), followed by
$g_\Gamma(\tau)-m_0=\Re\psi(a+i\tau/2)-\log\pi$, proves (177).

For (178), write $S=R_++R_-$ and $U=R_+/S$.  The beta--Gamma algebra says
that $S\sim\mathrm{Gamma}(2a,1)$, $U\sim\mathrm{Beta}(a,a)$, and $S,U$
are independent.  Also
$X=\tfrac12\log(U/(1-U))$ and
\[
 Y=\log\pi-\log S-\frac12\log(U(1-U)).
\]
For $a=1/4$, $U(1-U)=(4\cosh^2X)^{-1}$ and
$\mathbb E\log S=\psi(1/2)$, giving (178).  Positivity follows at
$x=0$ from $\psi(1/2)=-\gamma-2\log2$ and then from $\cosh x\ge1$.
The change of variables $u=e^{2x}/(1+e^{2x})$ in the beta density gives
(178a).
\end{proof}

Equation (177) is the exact Gamma--Tate normalization identity sought in
the preceding constructions: the Gamma multiplier and the constant channel
are respectively the ratio and scale matrix coefficients of one canonical
Mellin probability space.  It is stronger than merely checking the numerical
value of $m_0$.

It does not, by itself, prove (84c).  Passing from the positive Fourier
transform in (177) to the unweighted multiplier $m_0-g_\Gamma$ requires
division by $\phi_\Gamma$.  Stirling's formula gives
\begin{equation}
 \phi_\Gamma(\tau)^{-1/2}
 \asymp e^{\pi|\tau|/4}|\tau|^{1/4}
 \qquad(|\tau|\to\infty),                            \tag{179}
\end{equation}
so this deconvolution is unbounded on $L^2(\mathbb R)$ and on the general
Paley--Wiener image of compactly supported $L^2$ sources.  Consequently
(177) may be used as a weighted boundary identity, but treating its
deconvolution as a Hilbert-space contraction would require exponential
regularity absent from the ambient compact-window source.

In fact the two Tate constraints cannot supply that regularity.

\begin{proposition}[Gamma deconvolution domain no-go --- \PROVED]
Let $F\in L^2(I_T)$ and let $E_TF$ be its zero extension.  If
\begin{equation}
 \phi_\Gamma^{-1/2}\widehat{E_TF}\in L^2(\mathbb R), \tag{180}
\end{equation}
then $F=0$.  This remains true after restriction to the primitive space
$\mathcal P_T$.  Therefore no nonzero physical compact-window source can be
obtained by applying the unbounded Gamma deconvolution in (179), even after
the Tate moment conditions.
\end{proposition}

\begin{proof}
By Stirling's formula, (180) implies
$e^{b|\tau|}\widehat{E_TF}(\tau)\in L^2(\mathbb R)$ for every
$0<b<\pi/4$ (a polynomial factor is harmless after slightly decreasing
$b$).  Fourier inversion and Cauchy--Schwarz then extend $E_TF$ to a
holomorphic function on the strip $|\Im z|<b$.  Its boundary value on the
real axis vanishes on each open half-line outside $[-T,T]$.  The identity
theorem therefore gives $E_TF=0$.  The primitive space is a subspace of
$L^2(I_T)$, so the same conclusion holds there.
\end{proof}

Thus the weighted matrix coefficient (177) is a valid normalization theorem,
but deconvolution is not the missing G40 observation.  A successful bridge
must retain the Gamma ratio law inside its Hilbert metric and combine it with
the Euler/Tate channels before taking the physical compression.

This combined metric is exact in the absolutely convergent half-plane.

\begin{theorem}[Coherent completed Euler--Gamma metric --- \PROVED]
For $t>2$ put $r_p(t)=p^{-t/2}$ and evaluate the finite/infinite product
Poisson density on the scaling orbit $z_p(\tau)=p^{-i\tau}$:
\begin{equation}
 D_{E,t}(\tau):=prod_p
 \frac{1-r_p(t)^2}{|1-r_p(t)z_p(\tau)|^2}
 =\frac{|\zeta(\frac t2+i\tau)|^2}{\zeta(t)}.        \tag{181}
\end{equation}
Let
\begin{equation}
 \phi_{\Gamma,t}(\tau):=
 \frac{|\Gamma(\frac t4+\frac{i\tau}{2})|^2}
      {\Gamma(\frac t4)^2},
 \qquad
 c_t:=\pi^{-t/2}\Gamma\left(\frac t4\right)^2\zeta(t).
                                                               \tag{182}
\end{equation}
Then the completed metric has the exact factorization
\begin{equation}
 \left|\Lambda_{\mathbb Q}\left(\frac t2+i\tau\right)\right|^2
 =c_tD_{E,t}(\tau)\phi_{\Gamma,t}(\tau).             \tag{183}
\end{equation}
Thus the Euler Poisson boundary and the Gamma ratio boundary are not merely
analogous: their normalized coherent-orbit density is exactly the completed
Zeta metric, up to the scalar $c_t$.
\end{theorem}

\begin{proof}
Absolute convergence for $t>2$ permits multiplication of the Euler factors:
\[
 \prod_p(1-p^{-t})=\zeta(t)^{-1},\qquad
 \prod_p|1-p^{-t/2-i\tau}|^{-2}
 =|\zeta(t/2+i\tau)|^2,
\]
which proves (181).  Multiplication by (182) and the definition
$\Lambda_{\mathbb Q}(s)=\pi^{-s/2}\Gamma(s/2)\zeta(s)$ prove (183).
\end{proof}

The identity also locates the critical difficulty without hiding it in a
formal product.  As $t\downarrow1$, $c_t$ has the pole of $\zeta(t)$ while
the normalized Euler density loses its ordinary product-probability meaning;
only their product in (183) has the meromorphic continuation used in (119).
Consequently a critical Hilbert compression must renormalize the pair
$c_tD_{E,t}$ jointly.  Substituting the pointwise continuation of (183) for
such a renormalized boundary operator would be circular.

The required renormalization cannot be a vector limit in the original
$\ell^2$ completion.

\begin{theorem}[Euler coherent-orbit escape of mass --- \PROVED]
For $t>1$ let $e_t$ be the unit vector (131), and let
$W_\tau\delta_n=n^{-i\tau}\delta_n$ on $\ell^2(\mathbb N)$.  Then
\begin{equation}
 \left\langle W_\tau e_t,W_\sigma e_t\right\rangle
 =\frac{\zeta(t+i(\tau-\sigma))}{\zeta(t)}.           \tag{184}
\end{equation}
For every fixed $\tau$, $W_\tau e_t\rightharpoonup0$ as $t\downarrow1$,
although its norm is one.  Moreover
\begin{equation}
 \lim_{t\downarrow1}
 \left\langle W_\tau e_t,W_\sigma e_t\right\rangle
 =\begin{cases}1,&\tau=\sigma,\\0,&\tau\ne\sigma.
 \end{cases}                                          \tag{185}
\end{equation}
For $t>2$, the formal boundary-sum functional is absolutely defined on this
vector and satisfies
\begin{equation}
 \left|\sum_{n\ge1}(W_\tau e_t)_n\right|^2
 =\frac{|\zeta(\frac t2+i\tau)|^2}{\zeta(t)}
 =D_{E,t}(\tau).                                      \tag{186}
\end{equation}
Thus (181) is the squared boundary evaluation of the coherent vector, but
that evaluation and the vector itself have no ordinary critical limit.
The kernel (185) also cannot be realized by nonzero vectors indexed by all
$\tau\in\mathbb R$ in a separable Hilbert space.
\end{theorem}

\begin{proof}
Expanding in the standard basis gives
\[
 \langle W_\tau e_t,W_\sigma e_t\rangle
 =\zeta(t)^{-1}\sum_{n\ge1}n^{-t-i(\tau-\sigma)},
\]
which is (184).  Every fixed coordinate tends to zero because
$\zeta(t)^{-1/2}\to0$; boundedness of the norms and density of finitely
supported vectors imply weak convergence to zero.  If $\tau\ne\sigma$,
$\zeta(t+i(\tau-\sigma))$ has a finite limit while $\zeta(t)\to\infty$,
proving (185).  Equation (186) follows by summing the coordinates when
$t/2>1$.  Finally, a separable Hilbert space contains at most a countable
orthonormal set, whereas the limiting kernel in (185) indexes an uncountable
one.
\end{proof}

The gcd limit (134) uses the multiplicative isometries $V_n$, not the
continuous diagonal orbit $W_\tau$.  The two constructions therefore capture
different matrix coefficients of the same precritical vector $e_t$.  A valid
Poisson--Tate bridge must join them after removing the escaping polar
direction at the level of quadratic forms.  Identifying either limit with the
other without this quotient is an \AUDIT{} error.

There is nevertheless a canonical critical limit at the level of forms.

\begin{theorem}[Rigged critical Euler form --- \PROVED]
For $t>1$ set $k_t(u)=\zeta(t+iu)$.  If $f\in\mathcal S(\mathbb R)$ and
\begin{equation}
 \widehat f(x):=\int_{\mathbb R}f(\tau)e^{-i\tau x}\,d\tau,
\end{equation}
then
\begin{align}
 \mathfrak q_t(f)
 &:=\iint_{\mathbb R^2}f(\tau)\overline{f(\sigma)}
       k_t(\tau-\sigma)\,d\tau d\sigma                 \notag\\
 &=\sum_{n\ge1}n^{-t}|\widehat f(\log n)|^2\ge0.      \tag{187}
\end{align}
Moreover
\begin{equation}
 \lim_{t\downarrow1}\mathfrak q_t(f)
 =\mathfrak q_1(f)
 :=\sum_{n\ge1}\frac1n|\widehat f(\log n)|^2.         \tag{188}
\end{equation}
Hence $k_t\to\zeta(1+i\,\cdot)$ in the sense of positive tempered
distributions, where the boundary value on the right is the Fourier transform
of the positive locally finite measure
\begin{equation}
 \mu_{E,\mathrm{crit}}:=\sum_{n\ge1}\frac1n\delta_{\log n}. \tag{189}
\end{equation}
The completion of the form (188) is
$L^2(\mu_{E,\mathrm{crit}})$ (modulo the null space of the sampling map).
It is the canonical form-level renormalization of the escaping coherent
orbit.
\end{theorem}

\begin{proof}
For $t>1$, absolute convergence of the Dirichlet series and Fubini give
(187).  For every $N>1$, Schwartz decay gives
$|\widehat f(\log n)|\ll_N(1+\log n)^{-N}$, and
\[
 \sum_{n\ge2}\frac1{n(1+\log n)^{2N}}<\infty.
\]
Dominated convergence proves (188).  Equations (187)--(188) identify the
distributional Fourier transform and prove positivity.  The last assertion
is the standard completion of a sampled $L^2$ form.
\end{proof}

This critical form produces the arithmetic reference weights without an
Euler product or a sign assumption.  On the vector-valued space
$L^2(\mu_{E,\mathrm{crit}};L^2(\mathbb R))$, define for a fixed window
\begin{equation}
 (D_Tg)(\log n):=
 \mathbf1_{\{2\le n<e^{2T}\}}\mathbf1_{\{\Lambda(n)>0\}}
 (\Lambda(n)\sqrt n)^{1/2}g(\log n).                 \tag{190}
\end{equation}
Then, exactly,
\[
 \|D_Tg\|^2
 =\sum_{\substack{2\le n<e^{2T}\\\Lambda(n)>0}}
   \frac{\Lambda(n)}{\sqrt n}\|g(\log n)\|^2.
\]
Taking $g(\log n)=J_{\log n,-}F$ recovers the complete arithmetic part of
$\|\mathcal X_TF\|^2$.  Thus the arithmetic reference port is now obtained
as a diagonal observation of a source-defined critical Euler form, rather
than inserted as an unrelated direct sum.

What (187)--(190) do not yet provide is the opposite $J_{\log n,+}$ channel
or its global mixing with Gamma and the two Tate terminals.  That is precisely
the observation, rather than the reference Gram, which remains in G40.

The Euler critical form and the Gamma ratio law admit a positive joint
completion.

\begin{theorem}[Joint critical Euler--Gamma rigged boundary --- \PROVED]
For $t>1$ let
\begin{equation}
 \phi_{\Gamma,t}(u)=
 \frac{|\Gamma(\frac t4+\frac{iu}{2})|^2}
      {\Gamma(\frac t4)^2},
 \qquad
 K_t^{E\Gamma}(u):=\zeta(t+iu)\phi_{\Gamma,t}(u).     \tag{191}
\end{equation}
Then $K_t^{E\Gamma}$ is a positive-definite tempered distribution.  For
every $f\in\mathcal S(\mathbb R)$ its quadratic form has a finite critical
limit
\begin{equation}
 \mathfrak q_1^{E\Gamma}(f)
 =\int_{\mathbb R}|\widehat f(x)|^2
       d(\mu_{E,\mathrm{crit}}*\mu_{\Gamma,1})(x)\ge0, \tag{192}
\end{equation}
where $\mu_{\Gamma,1}$ is the probability law of the Gamma ratio $X$ in
(174).  Equivalently,
\begin{equation}
 K_t^{E\Gamma}\longrightarrow
 K_1^{E\Gamma}:=\zeta(1+i\,\cdot)\phi_{\Gamma,1}
 \quad\hbox{as positive tempered distributions}.     \tag{193}
\end{equation}
Consequently
\begin{equation}
 \mathscr H_{E\Gamma}^{\mathrm{crit}}
 :=L^2(\mu_{E,\mathrm{crit}}*\mu_{\Gamma,1})          \tag{194}
\end{equation}
is a canonical positive critical Euler--Gamma boundary constructed without
$\Delta_T$, the Weil sign, or a zero-free hypothesis.
\end{theorem}

\begin{proof}
The first factor in (191) is the Fourier transform of
$\sum_n n^{-t}\delta_{\log n}$ and the second is the characteristic function
of the Gamma ratio.  Their product is therefore the Fourier transform of the
convolution of two positive measures.  The density (178a), with $1/4$
replaced by $t/4$, has exponential tails uniformly for $t$ in a compact
right-neighbourhood of one.  Combined with the Schwartz decay of
$\widehat f$ and the summability argument in the proof of (188), this gives
an integrable majorant for the convolved sampling sum.  Dominated convergence
proves (192)--(193).  Positivity and the Hilbert completion (194) follow.
\end{proof}

This theorem constructs the positive target space required in the open
Tate-boundary statement (34) at the level of the global critical
Euler--Gamma form.  To prove (35), one must still construct the actual
intertwiner $\Phi_T$ into the windowed subspace of (194) and prove that its
compressed norm equals the Schur boundary form $J_T$.  The target is now
source-defined independently of that sign; the remaining equality is no
longer allowed to choose its codomain ad hoc.

The relation between this form boundary and the gcd boundary is exact, not
heuristic.

\begin{theorem}[Divisor-incidence realization of the gcd boundary --- \PROVED]
Let
\begin{equation}
 \mathscr E_N:=\ell^2\left(\{1,\ldots,N\},\frac1k\right),
 \qquad H_N:=\sum_{k\le N}\frac1k,
\end{equation}
and, for $n\le N$, define the truncated Zeta-incidence column
\begin{equation}
 z_n^{(N)}(k):=\sqrt n\,\mathbf1_{\{n\mid k\}}.
                                                               \tag{195}
\end{equation}
Then
\begin{equation}
 \frac1{H_N}\langle z_m^{(N)},z_n^{(N)}\rangle_{\mathscr E_N}
 =\frac{(m,n)}{\sqrt{mn}}
   \frac{H_{\lfloor N/[m,n]\rfloor}}{H_N},           \tag{196}
\end{equation}
where $[m,n]$ denotes the least common multiple.  Consequently, for fixed
$m,n$,
\begin{equation}
 \lim_{N\to\infty}\frac1{H_N}
 \langle z_m^{(N)},z_n^{(N)}\rangle
 =\frac{(m,n)}{\sqrt{mn}}=K(m,n).                    \tag{197}
\end{equation}
The inverse incidence relation is the exact finite M\"obius formula
\begin{equation}
 \mathbf1_{\{k=n\}}
 =\sum_{d\le N/n}\mu(d)\mathbf1_{\{nd\mid k\}}
 =\sum_{d\le N/n}\frac{\mu(d)}{\sqrt{nd}}
       z_{nd}^{(N)}(k).                               \tag{198}
\end{equation}
Thus the critical gcd GNS state is the harmonic-density renormalization of
the Zeta divisor-incidence columns in the rigged Euler boundary, and its
whitening is precisely M\"obius inversion.
\end{theorem}

\begin{proof}
The simultaneous support of the two columns consists of multiples of
$[m,n]$.  Hence
\[
 \langle z_m^{(N)},z_n^{(N)}\rangle
 =\sqrt{mn}\sum_{j\le N/[m,n]}\frac1{j[m,n]}
 =\frac{(m,n)}{\sqrt{mn}}H_{\lfloor N/[m,n]\rfloor},
\]
which proves (196).  Since $H_{\lfloor N/c\rfloor}/H_N\to1$ for fixed $c$,
(197) follows.  Finally, if $n\mid k$, the last sum in (198) is
$\sum_{d\mid k/n}\mu(d)$, equal to one only when $k=n$; if $n\nmid k$ it
is zero.  This proves (198).
\end{proof}

\begin{corollary}[Deformed incidence/GNS equivalence --- \PROVED]
For $t>1$ equip sequences with weight $k^{-t}$ and put
$z_{n,t}(k)=n^{t/2}\mathbf1_{\{n\mid k\}}$.  With
$H_M^{(t)}=\sum_{j\le M}j^{-t}$,
\begin{equation}
 \langle z_{m,t}^{(N)},z_{n,t}^{(N)}\rangle
 =\left(\frac{(m,n)}{\sqrt{mn}}\right)^t
   H_{\lfloor N/[m,n]\rfloor}^{(t)}.                \tag{199}
\end{equation}
After $N\to\infty$ and normalization by $\zeta(t)$,
\begin{equation}
 \frac1{\zeta(t)}\langle z_{m,t},z_{n,t}\rangle
 =\left(\frac{(m,n)}{\sqrt{mn}}\right)^t
 =\langle V_ne_t,V_me_t\rangle.                     \tag{200}
\end{equation}
Thus the cyclic spans of normalized incidence columns and of the shifted
Zeta vector $e_t$ are canonically isometric.  Their common critical limit is
the gcd GNS space (134)--(137).
If $A_t\delta_k=k^{t/2}\delta_k$, regarded as a unitary from the standard
finite $\ell^2$ space to the space with weight $k^{-t}$, then the Zeta
transport (27a) satisfies the exact identity
\begin{equation}
 z_{n,t}^{(N)}=A_t\mathscr Z_t\delta_n.               \tag{201}
\end{equation}
Hence the incidence model is unitarily the same finite Zeta transport used
in G39, before passing to either critical limit.
\end{corollary}

\begin{proof}
The calculation proving (196), with $k^{-1}$ replaced by $k^{-t}$, gives
(199).  Letting $N\to\infty$ gives the factor $\zeta(t)$; (133) gives the
last equality in (200).  Equality of all Gram matrices defines the canonical
isometry of the generated cyclic spans, and their pointwise Gram limit is
(134).  Finally, at $k=nd$ the $k$-th coordinate of
$A_t\mathscr Z_t\delta_n$ is
$(nd)^{t/2}d^{-t/2}=n^{t/2}$, and it vanishes when $n\nmid k$; this is
(201).
\end{proof}

Equations (195)--(198) close the previously open compatibility between the
finite Zeta gauge used in G39 and the global gcd GNS space: both arise from
one divisor-incidence system, and the normalizing scalar is forced by the
escaping harmonic mass.  In particular, the unorthogonalized incidence
columns are nested under $N\mapsto N'$, solving the nesting defect of the
coordinate square roots $K_T^{-1/2}$ noted after (141).

This does not yet identify the Gamma/Tate opposite observation.  It does,
however, place its arithmetic divergence (27e)--(27h), its reference Gram
(139)--(149), and the M\"obius normalization (150)--(153) in a single
source-derived boundary model.

There is a necessary affine component at the Gamma place.

\begin{theorem}[Linear/affine Gamma boundary separation --- \PROVED]
Let $\mathscr K_\Gamma=\bigoplus_{j\ge0}L^2(\nu_{a_j})$ and let
$\mathbf U_\tau=\bigoplus_jU_{j,\tau}$.  The vector-valued function
\begin{equation}
 b_\Gamma(\tau):=igoplus_{j\ge0}
 a_j^{-1/2}(U_{j,\tau}-I)\Omega_j                  \tag{202}
\end{equation}
is a continuous $1$-cocycle:
\begin{equation}
 b_\Gamma(\tau+\sigma)
 =b_\Gamma(\tau)+\mathbf U_\tau b_\Gamma(\sigma),
 \qquad
 \|b_\Gamma(\tau)\|^2=g_\Gamma(\tau).              \tag{203}
\end{equation}
In contrast, in the infinite tensor-product Gamma-ratio representation,
the coherent orbit $\Omega_\Gamma(\tau)$ satisfies
\begin{equation}
 \langle\Omega_\Gamma(\tau),\Omega_\Gamma(0)\rangle
 =\phi_\Gamma(\tau),
 \qquad
 \|\Omega_\Gamma(\tau)-\Omega_\Gamma(0)\|^2
 =2(1-\phi_\Gamma(\tau))\le2.                         \tag{204}
\end{equation}
There is no bounded operator $B$ such that
$B(\Omega_\Gamma(\tau)-\Omega_\Gamma(0))=b_\Gamma(\tau)$ for all $\tau$.
Thus the Gamma coherent metric and the Gamma energy feature are respectively
the linear and affine parts of the boundary and cannot be identified by a
bounded change of coordinates.
\end{theorem}

\begin{proof}
The cocycle identity follows from
$U_{\tau+\sigma}-I=(U_\tau-I)+U_\tau(U_\sigma-I)$.
Equations (170)--(171) show that (202) belongs to the direct sum and give its
norm in (203).  The tensor-product coefficient is (172), and the difference
norm gives (204).  Finally $g_\Gamma(\tau)\to\infty$ like $\log|\tau|$ by
the digamma asymptotic, whereas the norm in (204) is bounded.  A bounded
$B$ would contradict these two estimates.
\end{proof}

Consequently the source-defined positive target obtained so far must be
enlarged from (194) to an affine Hodge boundary carrying both
$\mathscr H_{E\Gamma}^{\mathrm{crit}}$ and the cocycle space
$\mathscr K_\Gamma$.  The Gamma channel of $\mathcal X_T$ lands in the
second component by (202)--(203); the scale vector (177) belongs to the first.
The missing G40 intertwiner is therefore a coupling between these two
components and the moving Tate terminal, not an isometry between either pair
alone.

The arithmetic component of that coupling has an exact lossless realization.

\begin{theorem}[Prime Julia colligation and M\"obius innovation --- \PROVED]
For a prime $p$ put $r=p^{-1/2}$, $s=(1-r^2)^{1/2}$, and
\begin{equation}
 \mathcal J_p:=
 \begin{pmatrix}r&s\\s&-r\end{pmatrix}.              \tag{205}
\end{equation}
Then $\mathcal J_p$ is a selfadjoint unitary colligation and its transfer
function is the Blaschke factor
\begin{equation}
 \Theta_p(z)=-r+zs(1-zr)^{-1}s
 =\frac{z-r}{1-rz},\qquad |\Theta_p(z)|=1\quad(|z|=1). \tag{206}
\end{equation}
With the M\"obius whitening multiplier
$b_p(z)=(1-rz)/s$, the prime-power innovations (166) obey
\begin{equation}
 e_{p,k}(z)=b_p(z)\Theta_p(z)z^{k-1}.                 \tag{207}
\end{equation}
For every finite set of primes $S$, the tensor product
$\bigotimes_{p\in S}\mathcal J_p$ is unitary and has transfer multiplier
$\Theta_S=\prod_{p\in S}\Theta_p$.  Hence the complete finite arithmetic
whitening/innovation system is a source-defined lossless scattering network.
\end{theorem}

\begin{proof}
The rows of (205) are orthonormal, so $\mathcal J_p^*\mathcal J_p=I$.
The standard transfer formula $D+zC(I-zA)^{-1}B$, with
$A=r$, $B=C=s$, $D=-r$, gives (206).  Its boundary modulus is one because
$|z-r|=|1-rz|$ on the unit circle.  Finally
\[
 b_p(z)\Theta_p(z)z^{k-1}
 =\frac{(z-r)z^{k-1}}s=e_{p,k}(z),
\]
which proves (207).  Finite tensor products of unitaries are unitary, and
scalar transfer functions multiply under the corresponding cascade.
\end{proof}

Together, (195)--(207) prove amplitude typing, nested Gram, M\"obius
normalization, prime-power innovation and positive conservation for the
arithmetic \emph{label} network.  The remaining physical cell must identify
its lossless output with the $J_{\log n,+}$ delay observation while coupling
the scale vector (177), the Gamma cocycle (202), and the moving two-moment
square (109).  Only after that joint affine Gamma--Tate termination is proved
does the network output equal the full factor $\mathcal Y_T$.

The intrinsic normalization of the Archimedean cell can now be closed.

\begin{theorem}[Exact affine Gamma--Tate termination identity --- \PROVED]
For $t>0$ set $a_j(t)=2j+t/2$ and
\begin{align}
 g_{\Gamma,t}(\tau)
 &:=\sum_{j\ge0}\frac{2\tau^2}
 {a_j(t)(a_j(t)^2+\tau^2)},                            \tag{208}\\
 c_\Gamma(t)&:=\pi^{-t/2}\Gamma\left(\frac t4\right)^2.
\end{align}
With $\phi_{\Gamma,t}$ from (191), one has the exact logarithmic metric
identity
\begin{equation}
 2\partial_t\log\!\left(c_\Gamma(t)\phi_{\Gamma,t}(\tau)\right)
 =g_{\Gamma,t}(\tau)+\psi\left(\frac t4\right)-\log\pi. \tag{209}
\end{equation}
At the critical endpoint this becomes
\begin{equation}
 \left.2\partial_t\log\!\left(
 c_\Gamma(t)\phi_{\Gamma,t}(\tau)\right)\right|_{t=1}
 =g_\Gamma(\tau)-m_0.                                 \tag{210}
\end{equation}
Thus the positive Gamma coherent metric, its affine energy cocycle and the
constant Tate scale channel form one differentiable Hilbert metric cell; the
normalization $m_0$ is forced by its scalar metric derivative.
\end{theorem}

\begin{proof}
For one factor,
\[
 \partial_t\log\frac{a_j(t)^2}{a_j(t)^2+\tau^2}
 =\frac1{a_j(t)}-\frac{a_j(t)}{a_j(t)^2+\tau^2}
 =\frac{\tau^2}{a_j(t)(a_j(t)^2+\tau^2)}.
\]
The series and its derivative converge locally uniformly, so summing gives
$2\partial_t\log\phi_{\Gamma,t}=g_{\Gamma,t}$.  Also
\[
 2\partial_t\log c_\Gamma(t)
 =-\log\pi+\psi(t/4).
\]
Adding proves (209).  At $t=1$, (171) identifies
$g_{\Gamma,1}=g_\Gamma$ and
$\psi(1/4)-\log\pi=-m_0$, proving (210).
\end{proof}

Equation (210) is the non-deconvolved form of (177).  It supplies the exact
normalization of the affine Gamma--Tate termination in the moving-metric
picture: the coherent metric must be transported, not cancelled.  The
remaining G40 identity is therefore the equality between the Kato-transported
physical compression of this cell, joined to the Julia arithmetic cascade,
and the opposite delay block $\mathcal Y_T$, including the already computed
second fundamental form (109).

The Euler cell has the complementary exact normalization.

\begin{theorem}[Exact coherent Euler--Gamma bulk normalization --- \PROVED]
For $t>2$, with $s=t/2+i\tau$, one has
\begin{align}
 2\partial_t\log\!\left(\zeta(t)D_{E,t}(\tau)\right)
 &=2\Re\frac{\zeta'}{\zeta}(s)                       \notag\\
 &=-2\sum_{n\ge2}\Lambda(n)n^{-t/2}
                  \cos(\tau\log n),                  \tag{211}\\
 2\partial_t\log\!\left(
 c_\Gamma(t)\zeta(t)D_{E,t}(\tau)
 \phi_{\Gamma,t}(\tau)\right)
 &=\Re\psi\left(\frac t4+\frac{i\tau}{2}\right)
   -\log\pi
   -2\sum_{n\ge2}\Lambda(n)n^{-t/2}
                  \cos(\tau\log n).                 \tag{212}
\end{align}
The product inside the logarithm in (212) is exactly
$|\Lambda_{\mathbb Q}(s)|^2$.  Its distributional continuation to $t=1$,
paired with a source supported in $I_T$, is
\begin{equation}
 g_\Gamma(\tau)-m_0
 -2\sum_{2\le n<e^{2T}}\frac{\Lambda(n)}{\sqrt n}
       \cos(\tau\log n)=a_T(\tau).                  \tag{213}
\end{equation}
Thus the entire Gamma, constant and prime-power bulk symbol of the
authoritative row-(d) operator is the covariant logarithmic derivative of the
single positive coherent Euler--Gamma metric, with no fitted scalar and no
omitted prime powers.
\end{theorem}

\begin{proof}
Equation (181) gives
$\zeta(t)D_{E,t}(\tau)=|\zeta(s)|^2$.  Differentiation yields the first line
of (211), and the absolutely convergent logarithmic Euler derivative gives
the series.  Adding (209), or differentiating (183) directly, proves (212).
The distributional continuation and compact-support localization were proved
in (93) and (117)--(119): a translation by $\log n\ge2T$ has zero quadratic
pairing with the zero-extended source, and the endpoint contact has null
overlap.  At $t=1$, (210) supplies $g_\Gamma-m_0$, proving (213).
\end{proof}

Equations (208)--(213) complete the normalization identity of the proposed
G40 bridge.  They do not assert that the derivative of the compressed metric
is positive: a differentiable family of positive metrics need not be
monotone.  The last load-bearing step is now solely the Kato/Julia compression
identity which must turn that covariant derivative, after removal of (109),
into the squared norm of the physical output residual.

The moving Gamma boundary embeddings themselves are explicit.

\begin{theorem}[Beta--Gamma Markov transport --- \PROVED]
Let $0<t<t'$, put $a=t/4$, $b=t'/4$, and let $X_t$ have the Gamma-ratio law
with characteristic function $\phi_{\Gamma,t}$.  If $B_+,B_-$ are independent
$\mathrm{Beta}(a,b-a)$ variables and
\begin{equation}
 Z_{t,t'}:=\frac12\log\frac{B_+}{B_-},                \tag{214}
\end{equation}
then there is a coupling with
\begin{equation}
 X_t=X_{t'}+Z_{t,t'},\qquad X_{t'}\perp Z_{t,t'},     \tag{215}
\end{equation}
and
\begin{equation}
 \mathbb E e^{i\tau Z_{t,t'}}
 =\frac{\phi_{\Gamma,t}(\tau)}{\phi_{\Gamma,t'}(\tau)}
 =\left|
 \frac{\Gamma(a+\frac{i\tau}{2})\Gamma(b)}
      {\Gamma(a)\Gamma(b+\frac{i\tau}{2})}
 \right|^2.                                          \tag{216}
\end{equation}
Consequently
\begin{equation}
 (\mathcal C_{t,t'}f)(x)
 :=\mathbb E\bigl[f(x+Z_{t,t'})\bigr]                \tag{217}
\end{equation}
defines a Markov contraction
$L^2(\mu_{\Gamma,t})\to L^2(\mu_{\Gamma,t'})$ under the coupling (215).
Its action on characters is multiplication by the explicit coefficient
(216).  These maps give a source-defined differentiable transport of the
Gamma boundary; no spectral sign or choice of $\Delta_T$ occurs.
\end{theorem}

\begin{proof}
Take independent Gamma variables $R_{a,\pm}$ of shape $a$ and
$G_{b-a,\pm}$ of shape $b-a$.  The beta--Gamma algebra gives
$R_{b,\pm}=R_{a,\pm}+G_{b-a,\pm}$ and
$B_\pm=R_{a,\pm}/R_{b,\pm}$, with $B_\pm$ independent of
$R_{b,\pm}$.  Taking half the difference of logarithms proves (215).
The beta Mellin formula
$\mathbb EB^z=\Gamma(a+z)\Gamma(b)/(\Gamma(a)\Gamma(b+z))$
proves (216).  Finally (217) is conditional expectation in the joint model,
and is therefore positivity preserving, unital and contractive on $L^2$.
\end{proof}

The infinitesimal logarithmic derivative of the character coefficient in
(216) is $-g_{\Gamma,t}(\tau)/2$, while the scalar metric derivative is the
remaining term in (209).  Hence (214)--(217) give the previously missing
boundary motion whose covariant generator is the affine Gamma--Tate cell.
What remains is to tensor this Markov transport with the finite Julia cascade
and prove that compression by the moving two-moment projector agrees with the
authoritative old/new decomposition.

The Euler transport exists as well, but with the opposite orientation.

\begin{theorem}[Oppositely oriented Euler and Gamma Markov transports
--- \PROVED]
Let $0<t<t'$ and, for every prime, put
$r_{p,t}=p^{-t/2}$ and $q_p=p^{-(t'-t)/2}$.  If
$Z_{p,t}$ has Poisson law $\mu_{r_{p,t}}$ and $Y_p$ independently has
Poisson law $\mu_{q_p}$, then
\begin{equation}
 Z_{p,t'}\ \buildrel d\over=\ Z_{p,t}Y_p,             \tag{218}
\end{equation}
because Poisson measures on the circle satisfy
$\mu_r*\mu_q=\mu_{rq}$.  Hence, on every finite prime window,
\begin{equation}
 (\mathcal C^E_{t,t'}f)(z)
 :=\mathbb E_Y f(zY)                                  \tag{219}
\end{equation}
is a Markov contraction
$L^2(\mu_{E,t'})\to L^2(\mu_{E,t})$.
The Gamma contraction (217), however, has direction
\begin{equation}
 \mathcal C^\Gamma_{t,t'}:
 L^2(\mu_{\Gamma,t})\longrightarrow
 L^2(\mu_{\Gamma,t'}).                               \tag{220}
\end{equation}
Thus the canonical Euler and Gamma conditional expectations point in
opposite directions.  There is no direct tensor-product Markov contraction
implementing both metric motions with the same time orientation.
\end{theorem}

\begin{proof}
The $k$-th Fourier coefficient of $\mu_r*\mu_q$ is
$r^{|k|}q^{|k|}=(rq)^{|k|}$, proving (218).  Jensen's inequality in the
coupling $Z_{t'}=Z_tY$ proves (219).  Equation (220) is (217).  Since the
domain of one canonical conditional expectation is the codomain of the other
with the time indices reversed, their tensor product is not a Markov map in
either common direction.
\end{proof}

This orientation theorem explains the role of the arithmetic Julia
colligation: its input/output swap is the source-defined lossless operation
available to reverse the Euler arrow before coupling it to (217).  The final
G40 bridge is now the concrete intertwining statement that this Julia-reversed
Euler transport, tensored with the beta--Gamma transport and then Kato
compressed by $P_t$, has output feature exactly $\mathcal Y_TP_T$ and defect
the terminal square (109) plus the old-compatible residual.

The infinitesimal Euler defects recover the authoritative prime-power weights
exactly.

\begin{theorem}[Prime-power Markov defect normalization --- \PROVED]
Let $\varepsilon=t'-t>0$.  On the $k$-th character of the $p$-th Poisson
factor, the Euler Markov contraction (219) acts by
\begin{equation}
 \mathcal C^E_{t,t'}z_p^k
 =p^{-k\varepsilon/2}z_p^k.                           \tag{221}
\end{equation}
Put $r_p=p^{-1/2}$.  The Julia defect of this scalar contraction, amplified
in an independent port for every prime power, has normalized squared energy
\begin{equation}
 d_{p,k}(\varepsilon)^2
 :=\frac{r_p^k}{k\varepsilon}
       \left(1-p^{-k\varepsilon}\right).             \tag{222}
\end{equation}
Then
\begin{equation}
 \lim_{\varepsilon\downarrow0}d_{p,k}(\varepsilon)^2
 =\frac{\log p}{p^{k/2}}
 =\frac{\Lambda(p^k)}{\sqrt{p^k}}.                  \tag{223}
\end{equation}
Thus every arithmetic coefficient in (83), including all prime powers, is
the positive infinitesimal defect rate of a source-defined Euler Markov/Julia
cell.  The factor $1/k$ is forced by the logarithmic Euler expansion and is
exactly cancelled by the character decay rate $k\log p$.
\end{theorem}

\begin{proof}
Equation (221) follows because the $k$-th moment of the independent Poisson
noise in (218) is $q_p^k=p^{-k\varepsilon/2}$.  A scalar contraction of
modulus $q_p^k$ has squared Julia defect $1-q_p^{2k}=1-p^{-k\varepsilon}$,
which gives (222) after the logarithmic Euler weight and the critical central
amplitude are inserted.  Finally
$(1-e^{-k\varepsilon\log p})/\varepsilon\to k\log p$,
proving (223).
\end{proof}

Equations (208)--(223) now derive all scalar normalizations of the proposed
bridge from positive metric motions: $g_\Gamma$, $m_0$ and every
$\Lambda(n)/\sqrt n$.  The remaining statement is not a scalar
normalization; it is the operator-level assertion that the amplified Julia
defect ports, after tensoring with the physical translated source and Kato
compressing by $P_t$, assemble with the plus orientation
$J_{\log n,+}$ and the old block required by (41p)--(41r).

The complete reference factor is the infinitesimal defect of these explicit
contractions.

\begin{theorem}[Markov-defect realization of $\mathcal X_T$ --- \PROVED]
Let
\begin{equation}
 \beta_\varepsilon(\tau):=
 \frac{\phi_{\Gamma,1}(\tau)}
      {\phi_{\Gamma,1+\varepsilon}(\tau)},
 \qquad
 d_{\Gamma,\varepsilon}(\tau)^2
 :=\frac{1-|\beta_\varepsilon(\tau)|^2}{\varepsilon}. \tag{224}
\end{equation}
For a prime power $n=p^k$ put
\begin{equation}
 d_{n,\varepsilon}^2
 :=\frac{p^{-k/2}}{k\varepsilon}
       (1-p^{-k\varepsilon}).                         \tag{225}
\end{equation}
On the form domain of $\mathcal X_T$, define the explicit defect analysis
map
\begin{equation}
 \mathcal D_{T,\varepsilon}F:=
 \left(
   \mathcal F^{-1}d_{\Gamma,\varepsilon}\widehat{E_TF},
   \{d_{n,\varepsilon}J_{\log n,-}F\}_{n\in\mathcal A_T}
 \right).                                             \tag{226}
\end{equation}
Then
\begin{equation}
 \lim_{\varepsilon\downarrow0}
 \|\mathcal D_{T,\varepsilon}F\|^2
 =\langle G_{\Gamma,T}F,F\rangle
  +\sum_{n\in\mathcal A_T}\frac{\Lambda(n)}{\sqrt n}
       \|J_{\log n,-}F\|^2
 =\|\mathcal X_TF\|^2.                               \tag{227}
\end{equation}
Every finite-$\varepsilon$ summand in (226) is a Julia defect of one of the
source-defined Markov contractions (217), (219).  Thus the full reference
Gram is a positive infinitesimal defect rate independently of row (d).
\end{theorem}

\begin{proof}
By (216) and (209),
\[
 -\log|\beta_\varepsilon(\tau)|^2
 =\int_1^{1+\varepsilon}g_{\Gamma,t}(\tau)\,dt.
\]
Hence $d_{\Gamma,\varepsilon}^2\to g_\Gamma$ pointwise.
Since $g_{\Gamma,t}\le g_{\Gamma,1}$ for $t\ge1$ and
$1-e^{-x}\le x$, the defect multipliers are dominated by $g_\Gamma$.
Dominated convergence on the Gamma form domain and (82) give the first
summand in (227).  Equations (222)--(223) give convergence of every
arithmetic coefficient; the contact set $\mathcal A_T$ is finite.  Summing
proves (227).  The defect interpretation is the definition of the Julia
defect operator $D_C=(I-C^*C)^{1/2}$ applied to the character eigenchannels.
\end{proof}

Therefore the input/reference half of the G40 network is now completely
constructed from positive contractions.  The remaining output construction
and identity are
\begin{equation}
 \boxed{\quad
 \operatorname*{form-lim}_{\varepsilon\downarrow0}
 \mathcal O_{T,\varepsilon}P_T
 =\mathcal Y_TP_T,
 \quad}                                                \tag{228}
\end{equation}
where $\mathcal O_{T,\varepsilon}$ must be constructed as the main output of
the coupled Julia dilations of the Euler contraction (219) and beta--Gamma
contraction (217), with the Euler ports reversed as required by (218)--(220),
and then transported by (112)--(116).  The admissible construction is thus
restricted to explicit source contractions, but its precise physical port
interconnection has not yet been proved.  Constructing it and proving (228),
with its orthogonal defect equal to (109) plus the old residual, would give
$\|\mathcal Y_TP_TF\|\le\|\mathcal X_TP_TF\|$ and close G40 without
circularity.

This probabilistic boundary realizes the positive metric
$\mathfrak M_t$ of (117) before critical continuation: the finite Euler
Poisson factors supply the Zeta modulus, while (172) supplies the Gamma
modulus.  The remaining constant/polar normalization and the two Tate
terminal motions are already explicit in (105)--(111).  The final G40 task
is therefore reduced to showing that the physical clamped-potential
observation is the conditional-expectation martingale associated with this
product boundary; then (157) and orthogonal innovations would give the
required unit contraction.

Thus the two ``explicit Tate terminals'' in (99) are no longer unspecified:
their canonical combined square is (109).  What remains in (104) is the
continuum Gamma--Poisson bulk comparison after subtracting this finite-rank
second-fundamental-form contribution.  The positivity of (109) alone cannot
fix the sign of that bulk term.

\section{Requirement-by-requirement completion audit}

\begin{center}
\begin{tabular}{p{0.20\textwidth}p{0.16\textwidth}p{0.54\textwidth}}
\hline
Requirement & Status & Authoritative evidence \\
\hline
Amplitude & \PROVED/\OPEN & (27g)--(27h) prove the finite divisibility
amplitude; identification with the continuum Hodge/old factor is not supplied.
\\
Normalization & \PROVED & (53)--(55) give the restricted Sonin polar
isometry without using row (d). \\
Gamma--Tate Gram & \PROVED & (81)--(84) give the exact source Gram and
balanced factorization; (138)--(149) construct a global gcd--Gamma reference
space, its Gram isometry and its positive midpoint conservation. \\
Typed observation & \PROVED/\OPEN & (84a)--(84c) define the canonical source
angular operator and (138)--(149) give the global cross-place type.  Its
contractive bound remains open.
\\
Continuous reduction & \PROVED & (90)--(94) identify the primitive space
with clamped potentials and give the exact Paley--Wiener defect kernel. \\
Global Poisson mixing & \PROVED/\OPEN & (96)--(98) give the full continuum
Euler divergence; its normalized physical Gram and boundary variation (104)
remain open. \\
Gamma normalization & \PROVED/\OPEN & (168)--(178a) give Gamma GNS and the
exact ratio--scale identity deriving $m_0-g_\Gamma$ from one Mellin law;
the unweighted deconvolution is excluded on every nonzero compact-window
source by (179)--(180).  Equations (181)--(183) identify the exact completed
Euler--Gamma coherent metric for $t>2$; its joint critical compression remains
open.  Equations (184)--(186) prove why it cannot be an ordinary vector
limit, while (187)--(194) construct its positive form-level and joint
Euler--Gamma replacement.  Equations (208)--(210) prove the exact
non-deconvolved Gamma--Tate metric termination. \\
Tate terminal motion & \PROVED & (105)--(111) give the exact moving
projection, minimum normal lift and positive even/odd terminal square. \\
Conservation & \PROVED/\OPEN & Krein/indefinite conservation and the defect
identity are exact; positive Hilbert conservation is equivalent to (84c). \\
Positivity & \OPEN & (84c) for every $T$ is exactly the Weil criterion and
hence equivalent to RH. \\
Old-factor compatibility & \PROVED/\OPEN & Restriction of (83) fixes the
authoritative $X_0,Y_0$ columns, while (41a)--(41j) construct the correct
signed old factor and capacity load.  Equations (195)--(201) identify the
finite Zeta gauge, gcd Gram and M\"obius normalization; equality with the
separate v83--v86 physical Hodge Gram (41) is still unproved. \\
Condition G / row (d) & \OPEN & The preceding two open Gram/passivity
requirements prevent a non-circular closure. \\
\hline
\end{tabular}
\end{center}

\section{Corrected final ledger after the type audit}

\subsection*{PROVED}
\begin{itemize}
\item The algebraic G39 divergence realization (27e)--(27h) in the finite
      divisibility model.
\item The restricted CCM polar isometry (52)--(55).
\item The unitary midpoint identity (44)--(47) as an abstract Hilbert-space
      theorem.
\item The typed observation criterion (57)--(61).
\item The order-mismatch no-go (62)--(63), excluding the row-(c) operator
      $K=\mathcal F^{-1}X\mathcal F$ as the Hodge energy feature.
\item The oriented-delay Kolmogorov factor and conservation identity
      (67)--(76) for its own positive reference Gram.
\item The exact Gamma delay formula (81)--(82) and the authoritative balanced
      factorization (83)--(84), including explicit factors
      $\mathcal X_T,\mathcal Y_T$.
\item The signed old-factor correction (41a)--(41c), the exact
      reference-harmonic capacity formula (41d)--(41i), and its equivalence
      (41j) with the return load in the supplied row (d).
\item The physical minimal amplitude current (41k)--(41m) and the complete
      joint-Gram criterion (41n)--(41o) for comparison with G39.
\item The finite-port/full-Hodge dimension obstruction: algebraic G39 cannot
      factor the infinite-dimensional physical old block.
\item The equivalence between the remaining sign and the unit Douglas
      contraction $C_T$.
\item The canonical source angular operator
      $\mathfrak C_T^{\rm src}$ and its exact defect pullback
      (84a)--(84b), constructed without a sign assumption.
\item The correctly typed observation theorem (86)--(87): normalized Sonin
      transport is isometric, whereas the physical observation must be
      contractive; an isometric observation would force $A_T=0$.
\item The equivalence of the all-window passivity estimate (84c), the Weil
      sign criterion, and RH.
\item The no-channelwise-contraction theorem (89): any valid $C_T$ must mix
      places and/or Gamma--Tate terminal channels.
\item The Tate-potential isomorphism and exact continuous kernel
      (90)--(94).
\item The differentiated global Poisson identity (96)--(98), the Gram
      homogeneity obstruction (100), and the exact Zeta metric variation
      (101)--(103).
\item The exact Tate boundary velocity and terminal energy (105)--(111).
\item The Kato primitive transport (112)--(116), the completed
      Euler--Gamma logarithmic metric (117)--(119), and the three-lines
      inverse-metric obstruction (120).
\item The exact residue Krein boundary and its fiber signatures
      (121)--(126), and the amenable unitarization theorem (127)--(128).
\item The critical gcd boundary state and explicit product-Poisson GNS
      realization (133)--(137).
\item The gcd--Gamma reference observation, global oriented midpoint
      colligation and exact source Gram (138)--(149).
\item M\"obius whitening, the centered Euler Fisher score, and the exact
      prime-power innovations (150)--(167).
\item The Archimedean Gamma GNS boundary (168)--(173), and the exact
      Gamma ratio--scale identity (174)--(178a), which derives both
      $g_\Gamma$ and $m_0$ from one canonical Mellin probability law.
\item The Gamma deconvolution domain no-go (179)--(180): neither compact
      support nor the two Tate moments permit removal of the Gamma
      characteristic factor as an $L^2$ operator.
\item The coherent completed Euler--Gamma factorization (181)--(183) in the
      absolutely convergent half-plane.
\item The Euler coherent-orbit escape theorem (184)--(186), separating the
      continuous scaling limit from the multiplicative gcd GNS limit.
\item The positive rigged critical Euler form, exact arithmetic diagonal
      observation, and joint Euler--Gamma boundary (187)--(194).
\item The divisor-incidence realization and deformed GNS equivalence
      (195)--(200), identifying the finite Zeta gauge, harmonic-mass
      renormalization, gcd boundary and M\"obius whitening in one model.
\item The linear/affine Gamma separation (202)--(204), identifying
      $g_\Gamma$ as an unbounded affine cocycle energy rather than a bounded
      coherent-orbit distance.
\item The prime Julia colligations and their exact M\"obius-innovation
      cascade (205)--(207), giving a lossless realization of the arithmetic
      label network.
\item The exact affine Gamma--Tate termination identity (208)--(210), deriving
      $g_\Gamma-m_0$ as the logarithmic derivative of one positive coherent
      metric cell.
\item The complete coherent Euler--Gamma bulk normalization (211)--(213),
      recovering the authoritative localized symbol $a_T$ with every prime
      power and the exact constant.
\item The beta--Gamma and oppositely oriented Euler Markov transports
      (214)--(220), supplying explicit moving boundary embeddings.
\item The prime-power Markov defect rates (221)--(223), deriving exactly
      every coefficient $\Lambda(p^k)/p^{k/2}$ from positive Julia defects.
\item The complete Markov-defect realization (224)--(227) of the physical
      reference factor $\mathcal X_T$.
\end{itemize}

\subsection*{OPEN}
\begin{itemize}
\item Compatibility of that feature with the fixed old Gram
      $A_0=\mathcal G_0^*\mathcal G_0$.
      Concretely, the three comparison identities (41p)--(41r).
\item Identification of the new source-derived gcd--Gamma reference space
      (138)--(149) with the independently defined CCM Sonin target, if that
      particular target is required; the source Gram/type isometry itself is
      now constructed.
\item The general-window Sonin lower-frame estimate
      \texttt{eq:framebound} from the supplied \texttt{main.tex}; its known
      comparison theorem covers only the prime-free window.
\item Construction of the normalized physical Poisson--Sonin Gram in (99)
      and proof of the boundary monotonicity/square identity (104).
\item Proof that the clamped-potential observation is the conditional
      expectation/martingale of the joint Gamma--Euler product boundary,
      while retaining the Gamma characteristic factor in the Hilbert metric;
      the direct deconvolution route is ruled out by (180).
      More precisely after (214)--(223): prove that the Julia-reversed Euler
      transport and beta--Gamma transport, amplified by the physical delay
      channels and Kato-compressed by $P_t$, have main output
      $\mathcal Y_TP_T$ and defect equal to (109) plus the old-compatible
      Schur residual.
\item Construction of the windowed intertwiner from the authoritative
      clamped-potential/old-new source into the positive joint critical space
      (194), proving its Gram is the Schur boundary form and making it
      compatible with the gcd GNS limit and moving Tate projection.
\item Construction of the Hilbert completion of the Poisson quotient and the
      uniform centered-scaling estimate (127), or an equivalent critical
      closed-range theorem in a source RKHS realization.
\item Construction, from source formulas, of the contraction
      $C_T\mathcal X_TP_T=\mathcal Y_TP_T$ with $\|C_T\|\le1$; equivalently,
      the unit-budget factorization of $Q_c$ through $D_0^{1/2}$.
      Equivalently, prove the global bound (84c) for the already constructed
      source angular operator.
\item Identification of the undefined symbol $\mathcal H_{5/4}$ with a
      precise Gamma form, if that notation is to be retained.
\item The positive factorization (43), hence G40, G41, G42, condition $G$, and
      row (d).
\end{itemize}

\subsection*{AUDIT}
\begin{itemize}
\item Equations (48)--(51) are only a candidate application of the midpoint
      theorem and are not a proof of G40.
\item Equation (56) is ill typed without (57), and it is false when its $K$ is
      interpreted as the row-(c) operator (62).
\item The auxiliary direct-sum G39 port cannot be the midpoint residual:
      Proposition (64)--(66) proves the required unitary orbit is impossible
      under that orthogonal embedding.
\item A full-space identification in (78) with the finite G39 port is ruled
      out by the dimension obstruction.  At most (78) can describe compatible
      finite compressions; it cannot close the continuum row (d).
\item The former identification
      $\mathcal B_T^*\mathcal B_T=\cL$ and
      $\|\mathcal B_TRv\|^2=S_M'$ was a redefinition, not a derivation from
      the supplied row-(d) operators; it has been withdrawn.
\item A direct sum of locally normalized Gamma/Euler isometries cannot supply
      the global contraction; theorem (89) gives an explicit primitive
      obstruction in every fixed arithmetic channel for large windows.
\item Positivity of the Sonin Hilbert--Schmidt Gram $M(f)$ does not identify
      it with the Gamma--Tate reference Gram and does not imply its open lower
      frame bound.
\item Three-lines convexity for the natural completed metric carries the
      wrong $|\Lambda_{\mathbb Q}|^2$ weight; cancelling it by division is a
      zero-free assumption.
\item The analogous attempt to cancel the Gamma characteristic factor in
      (177) is impossible on every nonzero compact-window source by (180),
      including on the primitive Tate subspace.
\item The Gamma ratio coherent orbit cannot replace the Gamma energy current:
      (202)--(204) prove that the former is bounded and the latter is an
      unbounded affine cocycle.
\item The gcd state limit and the continuous scaling-orbit limit cannot be
      identified as ordinary vectors: (185) proves escape of mass and an
      uncountable delta kernel.  They must be joined only after a polar/Tate
      quotient at form level.
\item The formal Krein completion and any square root of the desired Schur
      remainder remain equivalent to the unproved sign and cannot close G.
\end{itemize}

\end{document}
